\documentclass[10pt]{article}
\usepackage[T1]{fontenc}
\usepackage{lmodern}
\usepackage[a4paper,margin=1in]{geometry}
\usepackage{amsmath,amssymb,amsthm,mathtools}
\usepackage{microtype}
\usepackage[skip=0.45\baselineskip plus 2pt,indent=0pt]{parskip}
\usepackage{graphicx,booktabs,array,flafter}
\usepackage{etoolbox}
\usepackage{aliascnt}
\usepackage[colorlinks=true,linkcolor=blue,citecolor=blue,urlcolor=blue]{hyperref}
\usepackage[nameinlink,noabbrev,capitalize]{cleveref}
\numberwithin{equation}{section}
\newtheorem{theorem}{Theorem}[section]
\newaliascnt{lemma}{theorem}
\newtheorem{lemma}[lemma]{Lemma}
\aliascntresetthe{lemma}
\newaliascnt{proposition}{theorem}
\newtheorem{proposition}[proposition]{Proposition}
\aliascntresetthe{proposition}
\newaliascnt{corollary}{theorem}
\newtheorem{corollary}[corollary]{Corollary}
\aliascntresetthe{corollary}
\newaliascnt{claim}{theorem}
\newtheorem{claim}[claim]{Claim}
\aliascntresetthe{claim}
\theoremstyle{definition}
\newtheorem{definition}[theorem]{Definition}
\newtheorem{remark}[theorem]{Remark}
\newtheorem{problem}[theorem]{Problem}
\newcommand{\R}{\mathbb R}
\newcommand{\Z}{\mathbb Z}
\newcommand{\dcI}[1]{\{0,\ldots,#1-1\}}
\crefname{claim}{Claim}{Claims}
\Crefname{claim}{Claim}{Claims}
\crefname{section}{Section}{Sections}
\Crefname{section}{Section}{Sections}
\crefname{equation}{}{}
\Crefname{equation}{Equation}{Equations}
\setlength{\emergencystretch}{1.5em}
\allowdisplaybreaks[1]
\renewcommand{\topfraction}{0.85}
\renewcommand{\textfraction}{0.1}
\renewcommand{\floatpagefraction}{0.8}
\hypersetup{
 pdftitle={Unbounded GHZ dimension from fixed-order multiphoton sources},
 pdfauthor={Heng Li and Xizhi Liu},
 pdfsubject={Fixed source order, unbounded GHZ dimension, rigid hypergraphs}
}
\makeatletter
\patchcmd{\@maketitle}{\null\vskip 2em}{}{}{}
\makeatother
\title{\Large\bfseries Unbounded GHZ dimension from fixed-order multiphoton sources}
\author{
Heng~Li\thanks{School of Mathematics, Shandong University, Jinan, China, and Extremal Combinatorics and Probability Group (ECOPRO), Institute for Basic Science (IBS), Daejeon, South Korea. \href{mailto:heng.li@sdu.edu.cn}{\nolinkurl{heng.li@sdu.edu.cn}}.}
\and
Xizhi~Liu\thanks{School of Mathematical Sciences, University of Science and Technology of China, Hefei, China. \href{mailto:liuxizhi@ustc.edu.cn}{\nolinkurl{liuxizhi@ustc.edu.cn}}.}
}
\date{\today}
\begin{document}
\maketitle
\begin{abstract}
We study the largest local dimension of a Greenberger--Horne--Zeilinger (GHZ) state obtained by postselecting events with exactly one photon in each output path. When a source emits, it sends one photon with its internal label into every path in its support. The supports are distinct, and those with a given label partition the paths. For pair sources the maximum is 2 on every even number $n\ge6$ of paths. With triplet sources it grows without bound as $n$ increases through multiples of three, even when each source meets each of three fixed groups once. For fixed $r\ge3$ and every positive multiple $n$ of $r$, we construct a balanced $r$-partite network attaining local dimension at least $c_r n^{2(r-2)/7}$. In the tripartite construction, any choice of triples from different prescribed matchings that covers each path in the first two groups once misses a path in the third. A product of $r-2$ such constructions gives the bound for $r>3$. A local exchange argument gives the upper bound $O_r(n^{r-1}/\log n)$. On 9 paths the maximum is 13 for unrestricted triples and 6 with three fixed groups. When the $d\ge2$ prescribed perfect matchings are the only ones, emission probability $1/d$ at each independent coherent binary source maximizes the success probability for an equal-weight GHZ state.
All proofs were generated by GPT.
\end{abstract}

\section{Introduction}\label{sec:introduction}

On $n$ output paths, the Greenberger--Horne--Zeilinger (GHZ) state~\cite{GreenbergerHorneZeilinger1989,Lawrence2014} of local dimension $d$ is
\begin{equation}
 |\mathrm{GHZ}_{n,d}\rangle\coloneqq\frac{1}{\sqrt d}\sum_{i=1}^{d}|i\rangle^{\otimes n}.
 \label{eq:GHZ-definition}
\end{equation}
Across every nontrivial bipartition of the paths, this is a Schmidt decomposition with $d$ equal coefficients. We ask how large $d$ can become when each emission produces exactly $r$ photons and we retain events with one photon in every output path.

Krenn, Gu, and Zeilinger~\cite{KrennGuZeilinger2017} represented output paths by vertices and pair sources by edges. The sets of sources whose emissions put one photon in every path are precisely the perfect matchings. Gu, Chen, and Krenn~\cite{GuChenKrenn2020} extended this model to hypergraphs and gave coherent triplet-source constructions of local dimension 10 on 6 paths and 13 on 9 paths. Krenn et al.~\cite{KrennEtAl2026NoGo} asked whether the attainable dimension is unbounded for triplet sources.

Distinct sources have distinct spatial supports, and each emitted photon carries its source's internal label. For each label, the corresponding supports partition the paths. Bogdanov's graph theorem~\cite{Bogdanov2017} gives the maximum dimension for pair sources, namely 3 on 4 paths and 2 for even $n\ge6$. In the complex-weighted graph model, the ends of an edge may carry different labels. Distinct perfect matchings can then contribute to the same basis state and their amplitudes can interfere. For even $n$, Krenn et al.~\cite{KrennEtAl2026NoGo} ruled out $d\ge n$ for $n\ge4$, and Rajasekaran~\cite{Rajasekaran2026Lean} gave a Lean proof excluding $d\ge3$ for $n\ge6$.

For any fixed source order $r\ge3$, we obtain arbitrarily large local dimensions by increasing the number of paths.

\begin{theorem}\label{thm:ghz-growth}
Fix an integer $r\ge3$. For every positive multiple $n$ of $r$, there is a network of coherent $r$-photon sources whose normalized state after postselection on one photon per output path is $|\mathrm{GHZ}_{n,d}\rangle$, where $d\ge c_r n^{2(r-2)/7}$ and $c_r>0$ depends only on $r$. The paths form $r$ equal groups, with one photon per group from each emitting source. The sources can have the same emission probability, with each emission amplitude in phase with its vacuum amplitude.
\end{theorem}

For triples, a construction with $O(t^7)$ paths in each group has $t^2$ prescribed perfect matchings. A choice of triples from different matchings can cover every path in the first two groups exactly once. We assign the third path of each triple so that any such choice misses at least one path in the third group and covers another at least twice. For $r>3$, we combine $r-2$ such constructions on two common groups. Each factor has $t^2$ matching choices, and the $t^{2(r-2)}$ tuples give distinct GHZ branches.

In a rigid hypergraph, distinct perfect matchings share no edge, so any exchange that transforms one into another must replace every edge. This rules out short directed cycles in an auxiliary graph of allowed vertex replacements and yields $d=O_r(n^{r-1}/\log n)$. For triplet sources on 9 paths we prove that 13 is the largest attainable local dimension. If every triple must meet each of 3 fixed groups of size 3, the maximum falls to 6, and the extremal network is unique up to relabeling. The two 9-path upper bounds have elementary proofs.

Experiments have produced polarization GHZ states with 3 to 12 photons~\cite{BouwmeesterEtAl1999,ZhaoEtAl2004,WangEtAl2016,ZhongEtAl2018}, with one qubit encoded in the polarization of each photon. Erhard et al.~\cite{ErhardMalikKrennZeilinger2018} realized a 3-photon GHZ state of local dimension 3. The dimension in \cref{thm:ghz-growth} is the common Schmidt rank across all nontrivial cuts. Coherence between the GHZ branches is used in tests of local realism~\cite{GreenbergerHorneShimonyZeilinger1990,Lawrence2014} and quantum secret sharing~\cite{HilleryBuzekBerthiaume1999}.

\Cref{sec:model} identifies the accepted emission patterns with perfect matchings. \Cref{sec:asymptotic} constructs the networks, and \cref{app:construction} proves the lower bound for every part size. \Cref{sec:finite} proves the general upper bound and determines both 9-path maxima, with the unrestricted counting argument completed in \cref{app:finite}. \Cref{sec:resources} counts the source configurations and determines the optimal success probability for independent binary sources.

\section{The source model and rigid hypergraphs}\label{sec:model}

A source that emits into $r$ distinct paths specifies an $r$-element set of vertices. A collection of emissions places one photon in every path exactly when these sets partition the vertex set. We use this correspondence to compute the amplitude of each accepted string of internal labels.

\subsection{The combinatorial problem}

For an integer $r\ge2$, an \emph{$r$-uniform hypergraph}, or \emph{$r$-graph}, is a collection of $r$-subsets of a finite vertex set. All hypergraphs here are simple. We identify $H$ with its edge set and denote its vertex set by $V(H)$. For positive integers $q$, write $[q]\coloneqq\{1,\ldots,q\}$. A \emph{perfect matching} is a family of pairwise vertex-disjoint edges that covers $V(H)$. It has $m$ edges when $|V(H)|=rm$.

An $r$-graph is \emph{$d$-rigid} if its edge set is the disjoint union of $d$ perfect matchings $M_1,\ldots,M_d$ and these are its only perfect matchings. For $r\mid n$, let $f_r(n)$ be the largest $d$ for which there is a $d$-rigid $r$-graph on $n$ vertices. In a \href{https://mathoverflow.net/questions/347793/hypergraphs-with-only-disjoint-perfect-matchings}{2019 MathOverflow question}, Krenn asked for $f_r(rm)$ for all $r,m\ge2$. An $r$-graph is \emph{balanced $r$-partite} if its vertices can be partitioned into $r$ parts of equal size with each edge meeting every part once. Let $g_r(m)$ be the largest $d$ for a $d$-rigid $r$-graph of this kind with $m$ vertices in each part. Equivalently, in either class one may maximize the number of perfect matchings in a hypergraph whose perfect matchings are pairwise edge-disjoint. Deleting every edge that belongs to no perfect matching leaves precisely the union of those matchings.

In particular, $g_r(m)\le f_r(rm)$. Every vertex of a $d$-rigid hypergraph belongs to one edge of each $M_i$, so the hypergraph is $d$-regular. Consequently,
\[
 |H|=dm=\frac{dn}{r},
 \qquad
 f_r(n)\le\binom{n-1}{r-1},
 \qquad
 g_r(m)\le m^{r-1}.
\]
The union of any nonempty subfamily of the prescribed matchings has precisely the perfect matchings in that subfamily. Deleting edges creates no perfect matching, and each discarded matching loses all its edges.

For $n=r$, a simple $r$-graph has at most one edge, so $f_r(r)=1$. For $n=2r$, every edge has a unique possible partner in a perfect matching, namely its complement. The complete $r$-graph therefore attains $f_r(2r)=\frac12\binom{2r}{r}$. For example, $f_3(6)=10$, as used in the construction of Gu, Chen, and Krenn~\cite{GuChenKrenn2020}. With $r$ fixed pairs of paths, the $2^r$ edges that meet every pair once form $2^{r-1}$ complementary pairs. Thus $g_r(2)=2^{r-1}$, giving $g_3(2)=4$ on the same 6 paths. How large can $f_r(rm)$ and $g_r(m)$ be as $m$ grows with $r$ fixed?

Bogdanov~\cite{Bogdanov2017} proved that $K_4$ is the only simple graph with at least 3 perfect matchings in which all perfect matchings are pairwise edge-disjoint. An even cycle has exactly 2 perfect matchings, so $f_2(4)=3$ and $f_2(n)=2$ for every even $n\ge6$. Chandran and Gajjala~\cite{ChandranGajjala2026} studied edge colorings in which every perfect matching is monochromatic. They characterized the graphs for which the maximum number of colors represented by perfect matchings is 1 or 2.

Rigidity also implies uniqueness, up to a permutation of colors, of a proper $d$-edge-coloring. In a $d$-regular $r$-graph, each color class of a proper $d$-edge-coloring meets every vertex once and is therefore a perfect matching. In a $d$-rigid hypergraph these classes must be the prescribed $M_i$. A unique proper $d$-edge-coloring fixes the decomposition into $d$ perfect matchings. Rigidity also rules out any perfect matching outside that decomposition.

For fixed matching size $m\ge2$ and every $r\ge2$, Pohoata, Sauermann, and Zakharov~\cite[Theorem~4.3 of the preprint]{PSZ2025} construct a balanced $r$-partite $r$-graph on $rm$ vertices with at least $m^{r-D_m\sqrt r}$ prescribed perfect matchings, where $D_m$ depends only on $m$. Whenever $m$ disjoint edges are chosen from these matchings, all their matching indices must agree. This condition rules out shared edges and additional perfect matchings in their union. For fixed $r\ge3$, their rainbow matching bound~\cite[Theorem~1.2]{PSZ2025} concerns $m$ disjoint edges with pairwise distinct matching indices. In \cref{thm:growth} we also exclude mixed perfect matchings in which some indices repeat.

\subsection{A normalized source model}
\label{sec:optical-model}

Each source has a fixed spatial support and internal label. Down-conversion can produce photon pairs~\cite{HongMandel1985}, and path identity can make different emission alternatives indistinguishable~\cite{KrennHochrainerLahiriZeilinger2017}. The labels may be encoded in orbital angular momentum~\cite{MairEtAl2001}, in time bins defined by a train of pump pulses~\cite{DeRiedmattenEtAl2002}, or in frequency bins~\cite{OlislagerEtAl2010,KuesEtAl2017}. Distinct labels correspond to orthogonal states on each path.

Let $H$ be the union of edge-disjoint perfect matchings $M_1,\ldots,M_d$ on a set $V$ of $n=rm$ paths. Every photon emitted by a source in $M_i$ carries internal label $i\in[d]$. We also use $i$ as the color of that edge. In the figures and exported source lists, the internal labels run from $0$ to $d-1$. Since $M_i$ covers every path once, each mode $(v,i)$ belongs to exactly one source.

For $e\in M_i$, the source prepares a coherent superposition of the vacuum $|0_e\rangle$ and the state $|1_e\rangle$ with one photon in each of the $r$ modes $(v,i)$, $v\in e$. Let $p_e\in(0,1)$ be its emission probability and $\phi_e$ the phase of its emission amplitude relative to a common reference. The normalized joint state of the independent sources is
\begin{equation}
 |\Psi\rangle\coloneqq\bigotimes_{e\in H}\left(\sqrt{1-p_e}\,|0_e\rangle+e^{\mathrm{i}\phi_e}\sqrt{p_e}\,|1_e\rangle\right).
 \label{eq:normalized-source-state}
\end{equation}
Full coherence of the accepted superposition requires that neither the pump nor any unobserved mode retain information about which sources emitted. Idler-path alignment~\cite{ZouWangMandel1991}, many-particle path identity~\cite{Lahiri2018}, and temporal coherence~\cite{JhaEtAl2008} affect whether the emission alternatives can be distinguished. In energy--time interferometry, two path alternatives interfere in Franson's proposal~\cite{Franson1989} and three in the experiment of Thew et al.~\cite{ThewEtAl2004}.

Let $\Pi_1$ be the projection onto the subspace with exactly one photon in each spatial path, irrespective of its internal label. For a perfect matching $M$ of $H$, write $i_M(v)$ for the color of its edge through $v$. Summing the accepted terms of \cref{eq:normalized-source-state} over all perfect matchings $M$ of $H$ gives
\begin{align}
 \Pi_1|\Psi\rangle
 &=\sum_M B_M\bigotimes_{v\in V}|i_M(v)\rangle,\label{eq:matching-state}\\
 B_M&\coloneqq\prod_{e\in M}e^{\mathrm{i}\phi_e}\sqrt{p_e}\prod_{e\in H\setminus M}\sqrt{1-p_e}.\label{eq:matching-amplitude}
\end{align}
\begin{samepage}
\begin{proposition}\label{prop:source-rigidity}
Let $H$ be the union of $d$ edge-disjoint perfect matchings on $n$ paths. Associate with each edge $e\in M_i$ a source as in \cref{eq:normalized-source-state}, with $0<p_e<1$ and label $i$ on every path of $e$. The support of $\Pi_1|\Psi\rangle$ in the local label basis is $\{|i\rangle^{\otimes n}\colon i\in[d]\}$ if and only if $H$ is $d$-rigid. In this case, if the branch amplitudes $B_{M_i}$ have equal magnitude, the normalized postselected state is $|\mathrm{GHZ}_{n,d}\rangle$ up to a diagonal unitary on one output path.
\end{proposition}
\end{samepage}
\begin{proof}
At each path $v$, the local label identifies the unique incident edge of that color. Hence the label string $(i_M(v))_{v\in V}$ determines $M$. Each prescribed matching gives the same label on every path. Any other perfect matching gives a distinct, nonconstant label string. All coefficients in \cref{eq:matching-amplitude} are nonzero, so the constant strings exhaust the support exactly when $H$ is $d$-rigid. The normalized coefficients are then $B_{M_i}/(\sum_j|B_{M_j}|^2)^{1/2}$. A diagonal unitary on one path removes their phases. If their moduli agree, each coefficient becomes $1/\sqrt d$.
\end{proof}

In a more general source model, the observed label at a path need not identify the source that emitted into it. Different perfect matchings can then yield the same label string~\cite{KrennGuSoltesz2019,GuErhardZeilingerKrenn2019}. Their amplitudes contribute to the same basis state and can interfere, permitting other entangled states~\cite{GuChenZeilingerKrenn2019} and leading to a different extremal problem~\cite{ChandranGajjala2024}. Herzog et al.~\cite{HerzogEtAl1994} observed interference between two ways of creating a photon pair in the same crystal. Kysela et al.~\cite{KyselaEtAl2020} used path identity to generate 3-dimensional entanglement on two paths.

With ideal destructive photon counting, we retain records with one count in each output path. Heralding the state while keeping its photons would require a separate measurement. \Cref{fig:dimension-transition} compares the six-path examples and shows how a row and an exit index the prescribed matchings. \Cref{fig:comb} shows a choice of triples from two matchings that leaves a path in the third part uncovered.

\begin{figure}[htbp]
\centering
\includegraphics[width=\linewidth]{figures/main_dimension_transition.pdf}
\caption{Perfect matchings at fixed source order. (a) Maximum local dimension for pair sources and a lower bound for triplet sources in the model of \cref{sec:model}. (b) The grid shows the nine choices of row and exit for $t=3$ in the double comb of \cref{sec:asymptotic}. Each row and exit specify one projected matching. (c) Both matchings of a 6-cycle and three of the ten complementary pairs of triples in the six-path construction of Gu, Chen, and Krenn~\cite{GuChenKrenn2020}. The numbers on the vertices denote output paths.}

\label{fig:dimension-transition}
\end{figure}

\section{A polynomial lower bound at fixed source order}\label{sec:asymptotic}

The following bound gives an explicit value for the constant in \cref{thm:ghz-growth}.

\begin{theorem}\label{thm:growth}
For all integers $r\ge3$ and $m\ge1$,
\[
 f_r(rm)\ge g_r(m)\ge\left(\frac{m}{128(r-2)^3}\right)^{2(r-2)/7}.
\]
\end{theorem}

Taking $m=n/r$ in \cref{thm:growth} and applying \cref{prop:source-rigidity} gives \cref{thm:ghz-growth} with $c_r=[128r(r-2)^3]^{-2(r-2)/7}$. In particular, the same construction gives $g_3(n/3)\ge(n/384)^{2/7}$ for every positive multiple $n$ of 3.

One can also prescribe the number of branches exactly.

\begin{corollary}\label{cor:prescribed}
For integers $r\ge3$ and $d\ge2$, put $k=r-2$ and $m=\lceil128k^3d^{7/(2k)}\rceil$. There is a balanced $r$-partite $r$-graph on $n=rm$ vertices that is $d$-rigid and has $dm$ edges. In particular, $n=O_r(d^{7/(2(r-2))})$ and the number of source configurations is $O_r(d^{1+7/(2(r-2))})$.

\end{corollary}
\begin{proof}
By \cref{thm:growth}, this choice of $m$ gives at least $d$ prescribed perfect matchings. Retain any $d$ of them. Their union has exactly $d$ perfect matchings and $dm$ edges. For any $m'\ge m$, the same argument gives exactly $d$ perfect matchings and $dm'$ edges at part size $m'$.
\end{proof}

\Cref{tab:twelve-networks} gives a balanced tripartite construction on 39 paths for $d=12$.

\subsection{How the third part rules out additional perfect matchings}\label{sec:third-part}

A bipartite perfect matching can be represented as a directed cycle cover. Let the two parts $A$ and $B$ be copies of the same set. The arc $u\to v$ represents the pair $(u_A,v_B)$. Let $\mathcal P_1,\ldots,\mathcal P_q$ be all the cycle covers of the chosen digraph. For each occurrence of an arc in a cover, we choose a vertex in the third part $C$.

For an arc $e$ of cover $\mathcal P_i$, the \emph{tag} $\mathbf z_e+\mathbf v_i$ specifies the third output path of the corresponding triple. The anchor $\mathbf z_e$ depends only on the arc, and the translation $\mathbf v_i$ depends only on the cover. Each of the 3 photons carries the internal label $i$. We choose the vectors so that each prescribed matching covers $C$ exactly once.

In the 9-vertex example of \cref{fig:comb}, the projected covers are $\mathcal P_0=\{uv,vu,ww\}$ and $\mathcal P_1=\{uv,vw,wu\}$, where $uv$ denotes $u\to v$. The vertices $c_0,c_1,c_2$ are the three spatial paths in $C$. Take scalar translations $v_0=0$, $v_1=1$ and anchors $z_{uv}=0$, $z_{vu}=z_{wu}=1$, $z_{ww}=2$, $z_{vw}=-1$. In $M_i$, the triple that projects to $e$ uses the third path $c_{z_e+v_i}$, as shown in panels (c) and (d).

The projected covers share only the arc $u\to v$. In $M_0$, replacing the triple that projects to this arc by its counterpart in $M_1$ leaves every path in $A$ and $B$ covered once. In $C$ the new choice covers $c_1$ twice and misses $c_0$. The reverse switch in $M_1$ covers $c_0$ twice and misses $c_1$. These are the only choices of triples from both matchings that cover $A$ and $B$ exactly once. Neither covers $C$ exactly once, so $M_0$ and $M_1$ are the only perfect matchings.

\begin{figure}[htbp]
\centering
\includegraphics[width=\linewidth]{figures/main_selector_to_ghz.pdf}
\caption{A choice of triples that covers the first two parts can miss a path in the third. (a) The 2 projected cycle covers share the arc $u\to v$. (b) The triples that project to this arc use $c_0$ in $M_0$ and $c_1$ in $M_1$. Each triangle represents one source. Its color gives the internal label of the 3 photons, one in each path of the triple. Repeated vertex names in $A$ and $B$ denote copies in those parts. (c,d) For $i=0,1$, the sources of $M_i$ put one photon in every path and contribute the basis state $|i\rangle^{\otimes9}$ after projection. (e) Switching only the triple that projects to $u\to v$ covers $c_1$ twice and leaves $c_0$ uncovered.}
\label{fig:comb}
\end{figure}

For a choice of triples from several prescribed matchings, more than one tag may change. A perfect matching still has to use every element of $C$ once. We use the sum of its tags to rule out all such choices.

\begin{samepage}
\begin{lemma}\label{lem:exposed-translations}
Let $D$ be a digraph on $m$ vertices with loops allowed and no parallel arcs, and let $\mathcal P_1,\ldots,\mathcal P_q$ be all its distinct cycle covers. Fix a vector $\mathbf z_e$ for each arc and translations $\mathbf v_1,\ldots,\mathbf v_q$ in the same real vector space. Suppose that for each $i$ some linear functional takes a strictly smaller value at $\mathbf v_i$ than at every $\mathbf v_j$ with $j\ne i$. Let $C$ be an $m$-element subset of this space such that
\begin{equation}
 \{\mathbf z_e+\mathbf v_i\colon e\in\mathcal P_i\}=C\qquad(i\in[q]).
 \label{eq:con-selector-tags}
\end{equation}
On two copies $A,B$ of $V(D)$ and a third part $C$, insert $(u_A,v_B,\mathbf z_e+\mathbf v_i)$ for every arc $e=u\to v$ in $\mathcal P_i$. The resulting simple tripartite hypergraph is $q$-rigid.
\end{lemma}
\end{samepage}
\begin{proof}
By \cref{eq:con-selector-tags}, the $m$ triples from each cover form a perfect matching. Triples above the same arc from different covers have distinct tags, so these matchings are edge-disjoint. Any perfect matching projects to some $\mathcal P_i$. Both this matching and the prescribed matching above $\mathcal P_i$ use every element of $C$ once. Write $j(e)$ for the index of the chosen triple above $e$. Subtracting the two tag sums gives $\sum_{e\in\mathcal P_i}(\mathbf v_{j(e)}-\mathbf v_i)=\mathbf0$, since the anchors cancel. Apply a linear functional whose unique minimum on $\{\mathbf v_1,\ldots,\mathbf v_q\}$ occurs at $\mathbf v_i$. Each scalar summand is nonnegative and vanishes only if $j(e)=i$. Hence the matching consists of the triples from $\mathcal P_i$.
\end{proof}

For $t^2$ prescribed covers, put $\gamma_j\coloneqq\binom j2$ and take $\mathbf v_{ab}\coloneqq(a,b,\gamma_a+\gamma_b)$ for $0\le a,b<t$. The functional $\lambda_{ab}(x,y,z)\coloneqq2z-(2a-1)x-(2b-1)y$ satisfies
\begin{equation}
 \lambda_{ab}(\mathbf v_{a'b'}-\mathbf v_{ab})
 =(a'-a)^2+(b'-b)^2.
 \label{eq:main-exposure}
\end{equation}
Once each prescribed cover uses every tag in $C$ exactly once, the sum of squares in \cref{eq:main-exposure} rules out any mixed perfect matching. It remains to choose anchors satisfying \cref{eq:con-selector-tags} for every cover.

The projection for the growing family is the \emph{double comb}, whose routes are drawn in \cref{fig:doublecomb-routes}. For the route indexed by $(a,b)$, follow the upper chain of row $a$, leave at exit $b$, traverse the vertices of $K_{ab}$ in order, and return along the lower chain to the common vertex $s$. Every directed cycle of length at least 2 follows one of these routes through $s$. Since $s$ has no loop, a cycle cover uses exactly one such route. Loops cover the vertices outside the route. There are $t$ choices of row and $t$ choices within each row, giving $t^2$ projected matchings. With the tags assigned below, each of these covers extends to exactly one perfect matching of the tripartite hypergraph, giving $t^2$ GHZ branches.

\begin{figure}[htbp]
\centering
\includegraphics[width=\linewidth]{figures/main_double_comb_routes.pdf}
\caption{Three exits $b=0,1,2$ in one row of the double comb. Each panel shows the vertex $s$ twice. Each other dot represents a chain of vertices. Each colored route uses the common chain $F$ and the chain on its private block $K_b$. It visits $S_j$ for $j\le b$ and $P_\ell$ for $\ell\ge b$. A gray loop at a contracted chain represents one loop at each of its unused vertices.}
\label{fig:doublecomb-routes}
\end{figure}

Within a row, the exits whose routes use a given shared arc form the whole row, a prefix, or a suffix. We take the tags from a box of lattice points with selected points removed. For a loop at $x$, the exits at which its translated tag lies in a given removed rectangle form an interval. The interval counts show that enough anchors remain for the shared arcs. The tags still unused in each cover go to its private path. \Cref{app:construction} gives the details and the bound of $O(t^7)$ vertices per part.

For $m\ge2$, a cyclic extension gives $g_{r+1}(m)\ge2g_r(m)$. To see this, take a balanced $r$-partite rigid $r$-graph and index one part as $x_0,\ldots,x_{m-1}$. Add a part $y_0,\ldots,y_{m-1}$. Extend each edge through $x_i$ in two ways, by $y_i$ and by $y_{i+1}$, with indices modulo $m$. A perfect matching projects to an old perfect matching. If its edge through $x_i$ uses $y_{i+1}$, the edge through $x_{i+1}$ must use $y_{i+2}$. Thus every edge through $x_i$ uses $y_i$, or every edge through $x_i$ uses $y_{i+1}$. Each old matching has precisely these two extensions, and their edge sets are disjoint. The part size stays fixed as the total path count increases from $rm$ to $(r+1)m$.

\subsection{A product sharing two vertex parts}

Call a set $U$ of projection vertices an \emph{active support} if the loop at each vertex outside $U$ is its only incident arc. In a product, the factors may use the same first two parts if their active supports fit disjointly in a common index set.

\begin{theorem}\label{thm:support-product}
Let $k,m,q_1,\ldots,q_k$ be positive integers. For $f\in[k]$, let $H_f$ be a simple $q_f$-rigid tripartite $3$-graph with parts $A_f,B_f,C_f$ of size $m$. Identify $A_f$ and $B_f$ with two copies of a set $W_f$. Suppose there is a set $U_f\subseteq W_f$ such that every edge $(u_{A_f},v_{B_f},c)$ satisfies either $u,v\in U_f$ or $u=v\in W_f\setminus U_f$. If $\sum_{f\in[k]}|U_f|\le m$, there is a balanced $(k+2)$-partite hypergraph of uniformity $k+2$ and part size $m$ that is $\prod_{f\in[k]}q_f$-rigid.
\end{theorem}

Outside $U_f$, the first two coordinates must be equal. No projected pair crosses the boundary of $U_f$. Place the $U_f$ disjointly in a common index set, use two copies of that set for the first two parts, and retain a separate tag part $C_f$ for each factor. Choosing a perfect matching in each factor then gives one global perfect matching. Conversely, given a perfect matching of the product, keep only the tag coordinate for factor $f$ in each edge. In each edge whose first coordinate lies outside $U_f$, replace the second coordinate by the first. The resulting triples form a perfect matching of $H_f$. Rigidity forces every selected edge to have the same matching index in that factor. Applying this to all factors excludes additional global matchings. The complete construction and proof are in \cref{app:construction}.

For $k=r-2$ we use a box of lattice points whose three side lengths are of order $kT^2$, $kT^2$, and $kT^3$, respectively. This gives a tripartite factor with $(T+1)^2$ matchings. For this factor, we take the hub and all vertices whose loops are omitted in at least one prescribed cover as the active support. The bound $k|U|\le m$ lets us place $k$ copies on the two common parts, so the product has $(T+1)^{2k}$ matchings on parts of size $O(k^3T^7)$. This gives the exponent $2k/7$. At smaller part sizes a planar construction uses one row and satisfies the same support condition. The two support estimates are proved in \cref{app:rect-support,app:planar}.

\section{Upper bounds and sharp small cases}\label{sec:finite}

In a rigid hypergraph, distinct perfect matchings are edge-disjoint. We bound their number by considering which vertex replacements in one perfect matching can occur together. On 9 paths we determine the exact maximum both for unrestricted triplet sources and for sources that meet each of three fixed groups.

\subsection{An upper bound from partial exchanges}\label{sec:harmonic}

For ordinary graphs, Alon and Friedland~\cite{AlonFriedland2008} bounded the total number of perfect matchings in terms of the degree sequence. In a rigid hypergraph, an exchange that transforms one perfect matching into another must replace every edge. Set $\beta_1\coloneqq1$. For $m\ge2$, let $\beta_m$ be the smaller of $m/2$ and the harmonic sum $\sum_{j=1}^m 1/j$.

For an $r$-graph $H$, its $(r-1)$-shadow $\partial_{r-1}H$ consists of the $(r-1)$-sets contained in an edge. For $S\in\partial_{r-1}H$, its codegree $d_H(S)$ is the number of edges containing $S$.

\begin{theorem}\label{thm:dimension-upper}
Let $r\ge2$ and $m\ge1$ be integers, and let $H$ be a $d$-rigid $r$-graph on $rm$ vertices. Then
\begin{equation}
 d\beta_m \le\sum_{S\in\partial_{r-1}H}
 \left[1-\left(1-\frac1r\right)^{d_H(S)}\right]
 \le |\partial_{r-1}H|\le\binom{rm}{r-1}.
 \label{eq:harmonic-shadow}
\end{equation}
If $H$ is balanced $r$-partite with parts $V_1,\ldots,V_r$, then $d\beta_m\le\bigl|\{e\setminus V_1\colon e\in H\}\bigr|\le m^{r-1}$.
Consequently,
\begin{equation}
 f_r(rm)\le\left\lfloor\frac{1}{\beta_m}\binom{rm}{r-1}\right\rfloor,\qquad
 g_r(m)\le\left\lfloor\frac{m^{r-1}}{\beta_m}\right\rfloor.
 \label{eq:harmonic-extremal}
\end{equation}
\end{theorem}

\begin{proof}
If $m=1$, the only possible edge is the full vertex set. When it is present, $d=1$ and each of its $r$ subsets of size $r-1$ contributes $1/r$ to the sum in \cref{eq:harmonic-shadow}. In the partite case, deleting the vertex in $V_1$ from this edge leaves one projected set. Both bounds also hold for the empty hypergraph.

Suppose $m\ge2$ and fix a perfect matching $M$. For each $e\in M$ and $v\in e$, draw an arc $v\to w$ if $w\notin e$ and $(e\setminus\{v\})\cup\{w\}$ is an edge of $H$. Choose a set $Q$ consisting of one representative from each edge of $M$, and restrict this digraph to $Q$. Write $d_Q^+(v)$ for the outdegree of $v$ in the restriction.

A directed cycle of length less than $m$ would let us replace each representative on the cycle by its successor, producing a different perfect matching with at least one edge of $M$ unchanged. Hence the restricted digraph is acyclic or consists of a directed $m$-cycle. Any additional arc between vertices of such a cycle would create a shorter directed cycle. In the acyclic case a topological ordering $v_1,\ldots,v_m$ gives $d_Q^+(v_i)\le m-i$, so the sum below is at least $\sum_{j=1}^m 1/j$. On the directed cycle every outdegree is 1 and the sum is $m/2$. In either case,
\begin{equation}
 \sum_{v\in Q}\frac1{1+d_Q^+(v)}\ge\beta_m.
 \label{eq:harmonic-acyclic}
\end{equation}

Choose the representatives independently and uniformly within their matching edges. Fix $e\in M$ and $v\in e$, and put $c=d_H(e\setminus\{v\})$. For each other edge $F\in M$, let $k_F$ count the outneighbors of $v$ in $F$. Conditional on $v\in Q$, its outdegree $d_Q^+(v)$ is a sum of independent Bernoulli variables with probabilities $k_F/r$, where $0\le k_F\le r$ and $\sum_F k_F=c-1$. Hence
\begin{align*}
 \mathbb E\!\left[\frac1{1+d_Q^+(v)}\,\middle|\,v\in Q\right]
 &=\int_0^1\prod_F\left(1-\frac{k_F}{r}(1-z)\right)\,\mathrm dz\\
 &\le\int_0^1\left(1-\frac{1-z}{r}\right)^{c-1}\,\mathrm dz
 =\frac r c\left[1-\left(1-\frac1r\right)^c\right].
\end{align*}
The inequality uses $1-ku\le(1-u)^k$ for integers $0\le k\le r$ and $0\le u\le1/r$. Since $\mathbb P[v\in Q]=1/r$, taking expectations in \cref{eq:harmonic-acyclic} yields
\[
 \beta_m\le\sum_{e\in M}\sum_{S\in\binom e{r-1}}
 \frac{1-\bigl(1-\tfrac1r\bigr)^{d_H(S)}}{d_H(S)}.
\]
Now sum over the $d$ perfect matchings, which partition $H$. Each $(r-1)$-set $S$ in the shadow occurs exactly $d_H(S)$ times, so its denominator cancels. This proves \cref{eq:harmonic-shadow}.

For the partite bound, take $Q=V_1$. For $e\in M$, let $v$ be its unique vertex in $V_1$. All vertices completing $e\setminus\{v\}$ lie in $V_1$, so $1+d_Q^+(v)=d_H(e\setminus\{v\})$. Summing \cref{eq:harmonic-acyclic} over the $d$ matchings counts each distinct set of the form $e\setminus V_1$, $e\in H$, with total weight 1. There are at most $m^{r-1}$ such sets. Taking the maximum over the respective classes gives \cref{eq:harmonic-extremal}.
\end{proof}

For $m\ge5$ we have $\beta_m=\sum_{j=1}^m 1/j=\log m+O(1)$. For fixed $r$ and $n=rm\to\infty$, \cref{eq:harmonic-extremal} gives $f_r(n)\le(1/(r-1)!+o(1))\,n^{r-1}/\log n$ and $g_r(m)\le(1+o(1))m^{r-1}/\log m$. For triples this gives $f_3(n)=o(n^2)$, answering a question of \texttt{discrete\_things}~\cite{DiscreteThings2024}. For $f_3(n)$, the lower bound from \cref{thm:growth} is $c_3n^{2/7}$ and the upper bound is $(1/2+o(1))n^2/\log n$.

At $m=2,3,4$, the values of $\beta_m$ are $1,3/2,2$, respectively. In particular, $g_r(3)\le2\cdot3^{r-2}$ and $g_3(4)\le8$. The six-matching construction in the Supplemental Material~\cite[Theorem~S1.4]{LiLiuSupplement} applies at every part size $m\ge3$. Thus on 12 paths with 4 paths in each of 3 groups we have $6\le g_3(4)\le8$.

\subsection{Nine paths}

\begin{theorem}\label{thm:finite-exact}
Under the assumptions of \cref{sec:model}, the largest GHZ dimension on 9 paths is $f_3(9)=13$ when triples may have any spatial support, and $g_3(3)=6$ when every triple must meet each of three fixed groups. In the tripartite case the 6-rigid hypergraph attaining equality is unique up to relabeling within its parts.
\end{theorem}

Gu, Chen, and Krenn~\cite[Sec.~III\,B]{GuChenKrenn2020} reported that dimension 13 is attainable on 9 paths with triplet sources. In 2024, \texttt{discrete\_things}~\cite{DiscreteThings2024} listed 13 perfect matchings on 9 vertices whose union has no others. Below we give an algebraic construction over $\mathbb F_3$, show that it is nonisomorphic to the 2024 example, and prove $f_3(9)\le13$.

For unrestricted triplet sources, the values of $f_3$ at 3, 6, and 9 paths determine the minimum number of paths for each dimension up to 13. The minimum is 6 paths for dimensions $2$ through $10$ and 9 paths for dimensions $11$ through $13$. Indeed, $f_3(3)=1$, $f_3(6)=10$, and deleting whole matching classes realizes any smaller positive dimension on the same paths. At these minima the networks require $2d$ and $3d$ source configurations, respectively.

\subsubsection*{Three fixed groups of paths}

\Cref{thm:dimension-upper} gives $g_3(3)\le 3^2/\beta_3=6$. To attain this bound, take three disjoint copies of $\mathbb F_3$ as the vertex parts. Let $H$ consist of the triples $(a,b,c)$ with $a+b+c\in\{0,2\}$. For $h\in\{0,2\}$ and $s\in\mathbb F_3$, set
\begin{equation}
 M_{h,s}\coloneqq\{(a,a+s,a+h-s)\colon a\in\mathbb F_3\}.
 \label{eq:finite-six-matchings}
\end{equation}
These 6 perfect matchings partition the 18 edges of $H$, since an edge determines $h=a+b+c$ and $s=b-a$.

It remains to check that $H$ has no other perfect matching. In any perfect matching, the sum of the vertex coordinates is zero in $\mathbb F_3$. Its 3 edge sums lie in $\{0,2\}$, so all 3 must have the same value $h$. Write its edges as $(a,b(a),c(a))$, where $b$ and $c$ are permutations of $\mathbb F_3$. Every permutation of $\mathbb F_3$ has the form $a\mapsto\varepsilon a+s$ with $\varepsilon\in\{1,-1\}$. Write $b$ in this form. If $\varepsilon=-1$, the equation $a+b(a)+c(a)=h$ makes $c(a)=h-s$ constant, which is impossible. Thus $\varepsilon=1$, and the matching is one of those in \cref{eq:finite-six-matchings}. This proves $g_3(3)=6$.

Now let $H$ be a balanced tripartite $6$-rigid $3$-graph with 3 vertices per part. Take $Q=V_1$ in the proof of \cref{thm:dimension-upper}. Since $6\beta_3=9$, all 9 pairs in $V_2\times V_3$ occur, and each of the 6 matchings attains equality in \cref{eq:harmonic-acyclic}. An acyclic exchange digraph would contribute at least $1+1/2+1/3=11/6>3/2$, so each digraph is a directed triangle. Its outdegrees are all 1, giving each pair exactly 2 completions in $V_1$. The same argument applies with any part as $Q$. Thus the nine missing triples contain every cross-part pair exactly once and specify a Latin square of order 3. After relabeling, the first row and column read $0,1,2$. The middle entry must be 2, since a 0 there would put a second 2 in the last column. This forces the unique reduced square of order 3, as recorded by McKay and Wanless~\cite{McKayWanless2005}. The other 18 triples of $V_1\times V_2\times V_3$ form the unique extremal hypergraph up to relabeling within parts.

\subsubsection*{An algebraic construction with 13 matchings}

We now allow a source to meet any 3 of the 9 paths. Identify the vertex set with $\mathbb F_3^2$. For $\boldsymbol\sigma\in\mathbb F_3^2$, define
\[
 H_{\boldsymbol\sigma}\coloneqq\left\{\{\mathbf x,\mathbf y,\mathbf z\}\in\binom{\mathbb F_3^2}{3}\colon \mathbf x+\mathbf y+\mathbf z=\boldsymbol\sigma\right\}.
\]
Take $H\coloneqq\bigcup_{\boldsymbol\sigma\in\{0,1\}^2}H_{\boldsymbol\sigma}$. We determine every perfect matching of $H$ by first considering one sum class.

The triples in $H_{\mathbf0}$ are the 12 affine lines of $\mathbb F_3^2$. Two such lines are disjoint precisely when they are parallel. Thus the 4 parallel classes are exactly the perfect matchings of $H_{\mathbf0}$, and they partition its edges.

Fix $\boldsymbol\sigma\ne\mathbf0$. Choose $\mathbf u\notin\langle\boldsymbol\sigma\rangle$ and put $\mathbf v\coloneqq\boldsymbol\sigma-\mathbf u$. The vectors $\mathbf u,\mathbf v$ are linearly independent. Set $T\coloneqq\{\mathbf0,\mathbf u,\mathbf v\}$. Its 9 translates are distinct triples with sum $\boldsymbol\sigma$. Indeed, invariance under a nonzero translation would make $T$ an affine line, contrary to linear independence. There are $\binom93-12=72$ triples with nonzero sum, and invertible linear maps act transitively on the 8 nonzero sums. Consequently $H_{\boldsymbol\sigma}$ has exactly 9 edges, all of which are translates of $T$.

The difference set satisfies
\[
 (T-T)\setminus\{\mathbf0\}
 =\{\pm\mathbf u,\pm\mathbf v,\pm(\mathbf u-\mathbf v)\}
 =\mathbb F_3^2\setminus\{\mathbf0,\boldsymbol\sigma,-\boldsymbol\sigma\}.
\]
It follows that two distinct translates $T+\mathbf x$ and $T+\mathbf y$ are disjoint exactly when $\mathbf x-\mathbf y=\pm\boldsymbol\sigma$. A family of 3 translates forms a perfect matching precisely when its translation vectors constitute a coset of $\langle\boldsymbol\sigma\rangle$. Thus $H_{\boldsymbol\sigma}$ has 3 perfect matchings, partitioning its 9 edges.

Finally, let $M$ be any perfect matching of $H$, and let $\boldsymbol\sigma_1,\boldsymbol\sigma_2,\boldsymbol\sigma_3\in\{0,1\}^2$ be its edge sums. Their total is the sum of all vertices, which is $\mathbf0$. In each coordinate, 3 elements of $\{0,1\}$ sum to zero in $\mathbb F_3$ only when they are all equal. Hence $\boldsymbol\sigma_1=\boldsymbol\sigma_2=\boldsymbol\sigma_3$. Every perfect matching lies in one sum class. The construction therefore has exactly $4+3+3+3=13$ perfect matchings, and these matchings partition its 39 edges.

Represent $x,y\in\mathbb F_3$ by the integers $0,1,2$ and label the path $(x,y)$ by the integer $3x+y$. Give the three sources of matching $M_i$ in \cref{tab:finite-thirteen} the internal label $i$. These 39 triplet-source configurations give the 13 GHZ branches in the model of \cref{sec:model}.

To compare this hypergraph with the one listed by \texttt{discrete\_things}~\cite{DiscreteThings2024}, join each pair of vertices of codegree 4. In each case the resulting graph consists of 3 disjoint triangles. The number of these triangles whose vertex set is a hyperedge is 3 in \cref{tab:finite-thirteen} and 2 in the 2024 list, so the hypergraphs are nonisomorphic.

\begin{table}[htbp]
\centering
\caption{The 13 perfect matchings of the algebraic construction. A string such as $012$ denotes the triple $\{0,1,2\}$. The 13 indexed matchings are pairwise edge-disjoint.}
\label{tab:finite-thirteen}
\begin{tabular}{c@{\qquad}ccc@{\qquad\qquad}c@{\qquad}ccc}
\toprule
Index&\multicolumn{3}{c}{Source configurations}&Index&\multicolumn{3}{c}{Source configurations}\\
\midrule
0&012&345&678&7&037&148&256\\
1&036&147&258&8&046&157&238\\
2&048&156&237&9&058&136&247\\
3&057&138&246&10&013&268&457\\
4&015&267&348&11&025&178&346\\
5&024&168&357&12&067&124&358\\
6&078&123&456&&&&\\
\bottomrule
\end{tabular}
\end{table}


\Cref{fig:sourceevents} shows two prescribed emission patterns and a mixed pattern with the same total photon number. Of these three, only the prescribed patterns place one photon in every path. When all 13 accepted branches have the same amplitude, the normalized postselected state is $|\mathrm{GHZ}_{9,13}\rangle$.

\begin{figure}[htbp]
\centering
\includegraphics[width=\linewidth]{figures/main_source_events.pdf}
\caption{Source configurations and photon-number projection in the 13-matching construction. Squares denote emitting triplet sources, circles denote output paths, and colors specify the internal labels of the emitted photons. (a,b) After projection, the emission patterns $M_0$ and $M_4$ from \cref{tab:finite-thirteen} contribute the basis states $|0\rangle^{\otimes9}$ and $|4\rangle^{\otimes9}$, respectively. (c) The mixed event sends 2 photons to path $2$ and none to path $5$, so $\Pi_1$ sends the corresponding state to zero. The formula below the panels gives the equal-weight superposition of all 13 accepted branches.}
\label{fig:sourceevents}
\end{figure}

\subsubsection*{The unrestricted upper bound}

Suppose that a rigid $3$-graph on 9 vertices has at least 14 perfect matchings. Retaining any 14 and deleting the other edges gives a $14$-rigid hypergraph $H$ with 42 edges. We prove in \cref{app:finite} that no such $H$ exists, using the following three steps.

First, every pair has codegree at most 4. Suppose a pair $ab$ has codegree $k\ge5$. Put $W\coloneqq V(H)\setminus\{a,b\}$ and $X\coloneqq\{x\in W\colon abx\in H\}$. Form an auxiliary graph on all 35 triples of $W$ by joining two disjoint triples if the unique vertex of $W$ outside their union lies in $X$. Color the edge by that vertex. The triples of $H$ contained in $W$ induce exactly $k$ pairwise disjoint edges, one of each color in $X$. Counting edges across a cut excludes $k=6,7$. For $k=5$, we count disjoint pairs of 2-element subsets of $X$.

Second, some pair has codegree 4, since the average pair codegree is $3\cdot42/36=7/2$. Fix such a pair $ab$. For this choice of $a,b$, the auxiliary graph on all 35 triples of $W$ is the disjoint union of a star $K_{1,4}$ and 3 vertex--edge incidence graphs of $K_4$. Exactly $42-(14+14-4)=18$ triples of $H$ are disjoint from $\{a,b\}$. Using regularity, the codegree bound, and the four disjoint edges of distinct colors forced by rigidity, we determine the 18 triples up to relabeling.

Finally, each of $a$ and $b$ must belong to 10 edges that avoid the other vertex. Since every edge belongs to a unique perfect matching, the forced configuration allows at most 9 such edges for at least one of the two vertices. This contradiction proves $f_3(9)\le13$.

A construction on 18 paths gives $f_3(18)\ge13$. For every integer $m\ge2$, some network of unrestricted triplet sources on $3m$ paths has at least 10 GHZ branches. We extend the finite examples by grafting. If both hypergraphs are $d$-rigid, the first has $3m$ vertices, and the second is balanced tripartite with part size $q$, their graft is $d$-rigid on $3(m+q-1)$ vertices. The Supplemental Material~\cite{LiLiuSupplement} gives the initial networks and the proof.

\section{Source configurations and optimal emission probabilities}\label{sec:resources}

Fix a $d$-rigid network with $m=n/r$ sources in each prescribed perfect matching. An accepted event occurs when precisely the sources of one matching emit. In the independent binary-source model, $m$ and $d$ determine the maximum probability of preparing an equal-weight GHZ state.

\subsection{Counting paths, sources, and modes}

In a $d$-rigid $r$-graph on $n=rm$ vertices, each edge specifies one $r$-photon source configuration. The network uses $dm$ configurations in all and $nd$ orthogonal modes, one for each choice of a path and a matching label. Integrated source arrays~\cite{SpringEtAl2017} and spatial multiplexing~\cite{CollinsEtAl2013} have been demonstrated for heralded single photons. The $r$-photon configurations in the present model require a common phase reference.

For fixed $r$, \cref{cor:prescribed} constructs a network of dimension $d$ using polynomially many paths and source configurations. Conversely, every such realization satisfies $d=O_r(n^{r-1}/\log n)$ by \cref{thm:dimension-upper}, so
\[
 n=\Omega_r\!\left((d\log d)^{1/(r-1)}\right).
\]
Indeed, $d\le\binom{n-1}{r-1}\le n^{r-1}$ gives $\log n\ge(\log d)/(r-1)$. Substituting into $n^{r-1}=\Omega_r(d\log n)$ gives the claimed lower bound on the number of paths.

\subsection{Equal-weight branches in the binary-source model}

Suppose first that all sources in \cref{eq:normalized-source-state} have emission probability $p$, with the same phase for each emission amplitude relative to its vacuum amplitude. Each prescribed matching uses $m$ of the $dm$ sources. The matchings give mutually orthogonal output states, so their probabilities add to give the success probability
\begin{equation}
 P_{m,d}(p)=d\,p^m(1-p)^{m(d-1)}.
 \label{eq:binary-success}
\end{equation}
For $d\ge2$, differentiating its logarithm gives $m/p-m(d-1)/(1-p)$. This derivative vanishes only at $p=1/d$, is positive below this point, and is negative above it. Thus
\begin{equation}
 P_{m,d}^{\max}=d^{1-m}\left(1-\frac1d\right)^{m(d-1)}.
 \label{eq:binary-success-optimal}
\end{equation}
At the optimum, the expected number of emitting sources is $dm/d=m$, so the expected total photon number is $rm=n$. Even at this setting, a trial succeeds only if the emitting sources form one of the $d$ perfect matchings.

Allowing the sources to have different emission probabilities gives the same maximum when all branches must have equal weight.

\begin{proposition}\label{prop:balanced-pump}
Let $H$ be a $d$-rigid $r$-graph on $n=rm$ vertices, with $d\ge2$, and let its sources satisfy \cref{eq:normalized-source-state}. Among choices $0<p_e<1$ for which all accepted GHZ branches have equal weight, the maximum success probability is given by \cref{eq:binary-success-optimal}. This maximum is attained exactly when $p_e=1/d$ for every $e\in H$.
\end{proposition}
\begin{proof}
Set $x_e\coloneqq p_e/(1-p_e)$. By \cref{eq:matching-amplitude}, the probability that precisely the sources in $M_i$ emit is
\[
 |B_{M_i}|^2=\frac{\prod_{e\in M_i}x_e}{\prod_{e\in H}(1+x_e)}.
\]
Equal branch probabilities give the same geometric mean $z=(\prod_{e\in M_i}x_e)^{1/m}$ in every matching. Jensen's inequality for the strictly convex function $y\mapsto\log(1+e^y)$, applied to the numbers $\log x_e$ in $M_i$, gives $\prod_{e\in M_i}(1+x_e)\ge(1+z)^m$. Multiplying these inequalities over the $d$ matchings gives $\prod_{e\in H}(1+x_e)\ge(1+z)^{dm}$. Since $P=\sum_{i=1}^d |B_{M_i}|^2$ and each branch has numerator $z^m$, \cref{eq:binary-success,eq:binary-success-optimal} give
\[
 P\le\frac{d\,z^m}{(1+z)^{dm}}
   =P_{m,d}\!\left(\frac{z}{1+z}\right)\le P_{m,d}^{\max}.
\]
Equality in the first inequality holds exactly when $x_e=z$ for every edge. Equality in the second inequality holds exactly when $z/(1+z)=1/d$. Thus equality throughout is equivalent to $p_e=1/d$ for every $e\in H$.
\end{proof}

For fixed $r,n,d$, every $d$-rigid network uses $dn/r$ source configurations and, by \cref{prop:balanced-pump}, has the same maximum success probability for an equal-weight GHZ state in the independent binary-source model. \Cref{tab:twelve-networks} compares three explicit networks with $d=12$ under different incidence restrictions. Deleting one matching from the 9-path example gives the first. The second is balanced tripartite. The third is also balanced tripartite, and any two distinct source supports meet in at most one path. A hypergraph with this property is \emph{linear}.

\begin{table}[htbp]
\centering
\caption{These three triplet-source networks each attain local dimension 12. On $n$ paths, each matching uses $m=n/3$ of the $12m$ source configurations. The last column lists the maximum success probability per trial for equal branch weights under ideal transmission and detection, computed from \cref{eq:binary-success-optimal}.}\label{tab:twelve-networks}
\begin{tabular}{@{}lrrr@{}}
\toprule
Restrictions on source supports & Paths $n$ & Source configurations & $P_{m,12}^{\max}$\\
\midrule
Unrestricted        & 9   & 36  & $3.93\times10^{-4}$\\
Balanced tripartite & 39  & 156 & $4.43\times10^{-19}$\\
Linear balanced tripartite & 117 & 468 & $6.02\times10^{-58}$\\
\bottomrule
\end{tabular}
\end{table}

At dimension 13, a successful trial requires exactly 3 of the 39 sources to emit in the 9-path network and exactly 6 of the 78 in the 18-path network. By \cref{prop:balanced-pump}, their respective optima are $P_{3,13}^{\max}\simeq3.32\times10^{-4}$ and $P_{6,13}^{\max}\simeq8.46\times10^{-9}$. \Cref{fig:resources} plots \cref{eq:binary-success} for both networks.

\begin{figure}[htbp]
\centering
\includegraphics[width=0.77\linewidth]{figures/main_postselection.pdf}
\caption{Postselection probability $P_{m,d}(p)$ for four triplet-source networks with $n=3m$ paths. Each curve uses a common emission probability $p$ for all sources in that network. Dots mark $p=1/d$. Both dimension-13 curves peak at $p=1/13$. A successful trial requires 3 emissions on 9 paths and 6 on 18 paths.}\label{fig:resources}
\end{figure}

For fixed $r$ and $\alpha>0$, suppose $d=\Theta(n^\alpha)$. Substituting $m=n/r$ into \cref{eq:binary-success-optimal} gives
\[
 \log P_{m,d}^{\max} =(1-m)\log d+m(d-1)\log\!\left(1-\frac1d\right)
 =-\frac{\alpha}{r}\,n\log n+O(n).
\]
From the construction in \cref{thm:growth}, retain $d=\Theta(n^{2(r-2)/7})$ perfect matchings. In this family, $\log P_{m,d}^{\max}=-\frac{2(r-2)}{7r}n\log n+O(n)$. For triplet sources the coefficient of $-n\log n$ is $2/21$. These networks use $\Theta(n^{1+2(r-2)/7})$ source configurations for fixed $r$.

We can also allow a source to emit more than one photon into each of its assigned modes. Suppose source $e$ prepares the normalized pure state $\sum_{k\ge0}c_{e,k}|k\rangle^{\otimes r}$. In a lossless experiment with exact photon-number postselection, terms with $k\ge2$ are rejected because they place at least 2 photons in every path of $e$. The amplitude of an accepted matching is then given by \cref{eq:matching-amplitude} with $c_{e,1}$ in place of $e^{\mathrm{i}\phi_e}\sqrt{p_e}$ and $c_{e,0}$ in place of $\sqrt{1-p_e}$. For identical sources with vacuum probability $p_0$ and probability $p_1$ of emitting one photon per assigned mode, the success probability is $P=d p_1^m p_0^{m(d-1)}$.

Suppose each photon is transmitted and detected independently with the same probability $\eta$ in every mode. If $P$ is the ideal success probability, the joint probability of ideal acceptance and detection of all $n$ photons is $\eta^nP$. This is the observed success probability per trial when all losses are flagged and the corresponding trials are discarded, or when losses occur after ideal photon-number selection. With unflagged losses before selection, emissions of more than $n$ photons can also produce one count per path.

\section*{Acknowledgments}

The second author thanks Rishikesh Gajjala, whose 2024 question about the number of perfect matchings in rigid $3$-uniform hypergraphs motivated this work.

H.L. was supported by the National Natural Science Foundation of China (12501487), the China Scholarship Council, and the Institute for Basic Science (IBS-R029-C4). X.L. was supported by the Excellent Young Talents Program (Overseas) of the National Natural Science Foundation of China.

\section*{Data and code availability}

The Supplemental Material~\cite{LiLiuSupplement} gives the finite constructions and thresholds for the matching size. The archive \nolinkurl{supplementary_data.zip} contains source lists for those constructions, a program for exporting networks, and a table of lower bounds for $1\le m\le1000$. It also includes exact-arithmetic checks of the rectangle counts and support estimates in \cref{app:construction}. The source package contains two independent C++17 checks of the 9-path upper bound proved in \cref{app:finite}.

\appendix

\section{Construction at every part size}\label{app:construction}

Put $k\coloneqq r-2$ and $u\coloneqq(m/(128k^3))^{2/7}$. To prove \cref{thm:growth}, we seek at least $u^k$ perfect matchings on $m$ vertices per part. A single matching suffices for $u\le1$, and cyclic extension gives the bound for $1<u\le2$. For $2<u<224^2$ we use a planar factor with one row of routes. For $u\ge224^2$ we use a rectangular factor in which the row and the exit can be chosen independently. Both factors have exactly $m$ vertices per part and an active support $U$ with $k|U|\le m$, so \cref{thm:support-product} applies to $k$ copies. Throughout this appendix $[a,b]$ denotes the integer interval from $a$ to $b$, empty if $a>b$, and $[q]_0=\{0,\ldots,q-1\}$.

\subsection{The product on two common parts}

\begin{proof}[Proof of \cref{thm:support-product}]
Since $\sum_f|U_f|\le m$, we can relabel each $W_f$ by the same $m$-element set $\Omega$ so that the sets $U_f$ are disjoint. Use these labels in both of the first two parts of each factor, and keep the tag parts $C_1,\ldots,C_k$ disjoint. In the $i$th matching $M_{f,i}$ of $H_f$, where $1\le i\le q_f$, write the triple with first coordinate $u$ as $(u,p_{f,i}(u),c_{f,i}(u))$. Here $c_{f,i}\colon\Omega\to C_f$ is a bijection, and $p_{f,i}$ is a permutation of $\Omega$ that fixes every point outside $U_f$.

For a tuple $\mathbf i=(i_1,\ldots,i_k)$ of matching indices, let $p_{\mathbf i}$ agree with $p_{f,i_f}$ on each $U_f$ and fix all remaining points. This is a permutation, since each $p_{f,i_f}$ preserves $U_f$. Take two copies $A,B$ of $\Omega$ and form
\[
 M_{\mathbf i}\coloneqq\bigl\{
 \bigl(u_A,p_{\mathbf i}(u)_B,c_{1,i_1}(u),\ldots,c_{k,i_k}(u)\bigr)
 \colon u\in\Omega\bigr\}.
\]
Each $M_{\mathbf i}$ is a perfect matching because every coordinate map is a bijection. Let $H$ be their union.

Project an edge $(u_A,v_B,c_1,\ldots,c_k)$ to factor $f$ by taking $(u,v,c_f)$ if $u\in U_f$ and $(u,u,c_f)$ otherwise. An edge of $M_{\mathbf i}$ projects to the triple of $M_{f,i_f}$ at $u$. Thus an edge common to two prescribed global matchings would belong, after projection, to the corresponding matchings in every factor. Since the matchings in each factor are edge-disjoint, the two index tuples must agree.

Let $M$ be any perfect matching of $H$. Every edge whose first coordinate lies in $U_f$ also has its second coordinate in $U_f$. Hence the permutation from $A$ to $B$ induced by $M$ restricts to a permutation of $U_f$. Projecting to factor $f$ keeps this permutation on $U_f$ and uses the identity on its complement. The projection also covers $C_f$ once, so it is a perfect matching of $H_f$, say $M_{f,i_f}$. If an edge of $M$ comes from $M_{\mathbf j}$, its projection also belongs to $M_{f,j_f}$. Edge-disjointness gives $j_f=i_f$ for every $f$. Thus $M$ consists of the $m$ edges of $M_{\mathbf i}$.
\end{proof}

\subsection{A rectangular construction}\label{app:rectangular}

\begin{proposition}\label{prop:rectangular-factor}
Let $k\ge1$ and $m\ge128k^3\,224^7$ be integers, and put $T=\lfloor(m/(128k^3))^{1/7}\rfloor$. There is a balanced tripartite hypergraph of part size $m$ that is $(T+1)^2$-rigid and has an active support $U$ satisfying $k|U|\le m$.
\end{proposition}

Put $t=T+1$, $\gamma_b=\binom b2$, and $D=\gamma_T$. We use the translations $\mathbf v_{ib}=(i,b,\gamma_i+\gamma_b)$ of \cref{eq:main-exposure}, with $0\le i,b\le T$. Let
\[
 A=4kT^2,\qquad B=6kT^2,\qquad
 H=\left\lfloor\frac{m-1}{AB}\right\rfloor,\qquad \rho=m-1-ABH.
\]
Let $X$ consist of $[A]_0\times[B]_0\times[H]_0$ and any $\rho$ points in its next layer $[A]_0\times[B]_0\times\{H\}$. The two projection parts are copies of $X\cup\{s\}$, where $s$ is the hub in \cref{fig:doublecomb-routes}. The tag set $C$ will consist of $m$ points of the box $Y=[A+T]_0\times[B+T]_0\times[0,H+2D]$. We construct $C$ by removing points from $Y$. A loop at $u\in X$ receives the tag $u+\mathbf v_{ib}$ in the cycle cover with row $i$ and exit $b$ whenever that cover uses the loop.

\paragraph{Removed points and their owners.}
We remove rectangles in the planes with first coordinate $t,2t,\ldots$. We call these the primary levels. The table gives the second coordinates and heights for each of the three types of rectangle.
\[
\begin{array}{c@{\qquad}ll}
 &\text{second coordinate}&\text{height}\\
 \text{horizontal}&[(2a+1)T,(2a+2)T-1]&[2D,H-1]\\
 \text{vertical}&[T,B-1]&[(2a+2)D,(2a+3)D-1]\\
 \text{full}&[T,B-1]&[2D,H-1].
\end{array}
\]
Here $0\le a<L_h=B/(2T)$ on a horizontal level and $0\le a<L_v=\lfloor(H-D)/(2D)\rfloor$ on a vertical level. We will also use at most one shorter rectangle of the same width as the full type. All primary levels have first coordinate at most $A-1$, so the inverse translates of their rectangles lie in the full layers of $X$.

In row $i$, a point $x\in X$ is an \emph{owner} of a primary rectangle $\mathcal R$ if $x+\mathbf v_{ib}\in\mathcal R$ for some exit $0\le b\le T$. The first coordinate of an owner modulo $t$ determines $i$. The inverse translates of rectangles on different primary levels lie in disjoint intervals of first coordinates. On a given level the unused blocks separate the inverse translates of distinct rectangles. Thus each owner is assigned to a unique rectangle and row.

For an owner $x$ of $\mathcal R$ in row $i$, the exits $b$ for which $x+\mathbf v_{ib}\in\mathcal R$ form a nonempty interval. If this interval is a proper prefix $[0,\ell]$ or suffix $[j,T]$, leave the loop at $x$ out of precisely the cycle covers at those exits. In every other case, leave it out of every cycle cover in the row. If the interval is internal, say $[a,b']$, the point $x$ can serve as a \emph{residual anchor} for a shared arc on each of $[0,a-1]$ and $[b'+1,T]$. We count its uses separately for the two patterns. The net demand of a prefix or suffix is the number of owners assigned that pattern minus the number of residual anchors available for it.

\begin{lemma}\label{lem:rectangle-owners}
Fix a row $i$ and a primary rectangle $\mathcal R$ of width $W\ge T$ and height $J\ge D$. Group the owners by the interval of exits at which their translates lie in $\mathcal R$. For $0\le\ell<T$, each of the patterns $[0,\ell]$ and $[\ell+1,T]$ has net demand $J-D+(W-T+1)\ell$. The number of owners assigned the whole row is $(W-T)(J-D)+(T-1)D$. The rectangle has $(W+T)J+(W-1)D$ owners in row $i$.
\end{lemma}
\begin{proof}
For a fixed $x$, the exits $b$ for which $x+\mathbf v_{ib}$ lies in $\mathcal R$ form the intersection of two intervals, one for each coordinate of the rectangle. If an owner's interval is $[0,\ell]$, either the horizontal interval ends at $\ell$ and the vertical one contains $[0,\ell]$, or the vertical interval ends at $\ell$ and the horizontal one contains $[0,\ell+1]$. For an internal interval starting at $\ell+1$, a suffix in one coordinate overlaps a prefix in the other away from both endpoints of $[0,T]$. Counting the owners in each case gives
\[
\begin{aligned}
 \#\{\text{owners with interval }[0,\ell]\}
   &=J-\gamma_\ell+\ell(W-\ell-1),\\
 \#\{\text{owners with internal interval starting at }\ell+1\}
   &=D-\gamma_{\ell+1}+\ell(T-\ell-1).
\end{aligned}
\]
Since $\gamma_{\ell+1}-\gamma_\ell=\ell$, their difference is $J-D+(W-T+1)\ell$. For each exit $b$, exactly $WJ$ points $x\in X$ satisfy $x+\mathbf v_{ib}\in\mathcal R$. Thus as many owners have intervals ending at $\ell$ as have intervals starting at $\ell+1$. Subtract the counts of internal intervals at both endpoints from this equality. The two sides are then the net demands for $[0,\ell]$ and $[\ell+1,T]$.

There are $(W-T)(J-D)$ owners with a full interval. The other owners assigned the whole row have internal intervals, formed where a prefix in one coordinate overlaps a suffix in the other strictly inside $[0,T]$. Counting these overlaps gives $\sum_{b=0}^{T-1}\gamma_b+\sum_{a=1}^{T}(D-\gamma_a)=(T-1)D$ such owners.

To count all owners, start with the $WJ$ points at exit 0. For $1\le b\le T$, any point of the new inverse translate that lies in an earlier translate also lies in the translate for $b-1$, by monotonicity in both coordinates. The height changes by $\gamma_b-\gamma_{b-1}=b-1$, so the new translate adds $WJ-(W-1)(J-b+1)=J+(W-1)(b-1)$ points. Summing these increments and using $\sum_{b=1}^T(b-1)=D$ gives the stated total.
\end{proof}

\paragraph{The shared arcs.}
In row $i$, place the owners whose loops are omitted on the full row, on the prefix $[0,\ell]$, or on the suffix $[j,T]$ on chains $F_i$, $P_{i,\ell}$, and $S_{i,j}$, respectively. Fix an order on the vertices of each chain. Each choice of row $i$ and exit $b$ also has a private block $K_{ib}$ of vertices, which the route traverses in a fixed order. Use the double comb of \cref{fig:doublecomb-routes} with routes
\[
 s,F_i,S_{i,1},\ldots,S_{i,b},K_{ib},P_{i,b},\ldots,P_{i,T-1},s
 \qquad(0\le i,b\le T).
\]
The suffix list is empty at $b=0$ and the prefix list at $b=T$. Insert all arcs along these routes and a loop at every point of $X$. Removing $s$ and the loops leaves an acyclic digraph. A cycle through $s$ chooses one row, passes through one private block and returns to $s$. Hence each route, together with the loops off it, is a cycle cover and these are all the covers. The shared arcs entering or lying inside a full-row or suffix chain are used at exactly the exits assigned to that chain. For a prefix chain the same holds for its internal and outgoing shared arcs. Thus a chain of $a$ owner vertices requires $a$ shared-arc anchors of its pattern. The arc entering $K_{ib}$ and the arc leaving it both belong to the private path.

Fix $i$. Each pattern has two sets of external anchors, listed below.
\[
\begin{array}{ll}
[0,T]&([-i,-1]\cup[A,A+T-i-1])\times[B]_0\times[H]_0,\\
     &[A]_0\times[B]_0\times([-\gamma_i,-1]\cup[H+1,H+D-\gamma_i]),\\[2pt]
[0,\ell]&[A]_0\times\{B+T-\ell-1\}\times[H]_0,\\
     &[A]_0\times[B]_0\times[H+1+2D-\gamma_i-\gamma_{\ell+1},H+2D-\gamma_i-\gamma_\ell],\\[2pt]
[j,T]&[A]_0\times\{-j\}\times[H]_0,\\
     &[A]_0\times[B]_0\times[-\gamma_i-\gamma_j,-\gamma_i-\gamma_{j-1}-1].
\end{array}
\]
For a fixed row, all the listed anchor sets are pairwise disjoint and lie outside $X$. The upper vertical intervals begin above $H$, so they also avoid the partial layer. At every exit in its pattern, each anchor translates to a point of $Y$. For an anchor in a horizontal piece, either the first coordinate of its translate lies outside $[T,A-1]$ or its second coordinate lies outside $[T,B-1]$. For an anchor in a vertical piece, its translated height lies below $2D$ or above $H$. We remove points only inside $[T,A-1]\times[T,B-1]\times[2D,H-1]$, so all these tags remain available.

\paragraph{Choosing the rectangles.}
Let $Q=B-2T+1$. For the moment allow real multiplicities $x,y,z$ for the horizontal, vertical, and full types, respectively. For the prefix $[0,\ell]$ and the suffix $[\ell+1,T]$, the external supply for each pattern is $AH+AB\ell$. The net demand is $x(H-3D+\ell)+yQ\ell+z(H-3D+Q\ell)$. For the full row the supply is $TBH+ABD$, and the demand is $(x+y+z)D(T-1)+z(B-2T)(H-3D)$. For the proper prefixes and suffixes, match the constant terms and the coefficients of $\ell$ in the demand and supply formulas. For the full row, equate total demand and total supply. Solving gives
\begin{equation}\label{eq:rect-calibration}
\begin{gathered}
 S=\frac{AH}{H-3D},\qquad x=S-z,\qquad y=\frac{AB-x}{Q}-z,\\[2pt]
 z=\frac{TBH+ABD-D(T-1)\bigl(S+\tfrac{AB-S}{Q}\bigr)}
 {(B-2T)\bigl(H-3D-\tfrac{D(T-1)}{Q}\bigr)}.
\end{gathered}
\end{equation}
Take $n_h=\lfloor x\rfloor$ horizontal rectangles, $n_v=\lfloor y\rfloor-1$ vertical rectangles, and $n_f=\lfloor z\rfloor$ rectangles of the full type, filling each horizontal or vertical level before starting another. Set $\alpha=z-n_f$ and $L=\lfloor\alpha(H-2D)\rfloor$. If $L\ge D$, add the shorter rectangle $[T,B-1]\times[2D,2D+L-1]$ at a new level. The resulting number of primary levels is
\begin{equation}\label{eq:rect-level-count}
 U_0=\left\lceil\frac{n_h}{L_h}\right\rceil
     +\left\lceil\frac{n_v}{L_v}\right\rceil+n_f+\mathbf 1_{\{L\ge D\}}.
\end{equation}

These integer choices leave enough anchors for every pattern. If the shorter rectangle is present, it adds $L-D+Q\ell$ to the net demand for each of $[0,\ell]$ and $[\ell+1,T]$. This is at most $\alpha(H-3D+Q\ell)+(1-\alpha)Q\ell$. Omitting one vertical rectangle saves $Q\ell$. The shorter rectangle's full-row demand minus $\alpha$ times that of a full rectangle is at most $-(1-\alpha)D(B-3T+1)\le0$. If the shorter rectangle is absent, rounding cannot increase demand. In \cref{app:rect-support} we verify $x>1$, $y>2$ and $z>T-1$, so all three types occur. Since $H>3D$ and $B>2T$, a horizontal rectangle gives owners of every proper pattern and a full rectangle gives full-row owners.

\paragraph{Completing the tag set.}
For a primary rectangle of width $W$ and height $J$, sum the net demands for the prefixes $[0,\ell]$ over $0\le\ell<T$ and add the full-row demand. The formulas in \cref{lem:rectangle-owners} and $\sum_{\ell=0}^{T-1}\ell=D$ give exactly its area $WJ$. At the real multiplicities in \cref{eq:rect-calibration}, each demand equals its external supply. The total primary area before rounding is therefore
\[
 \sum_{\ell=0}^{T-1}(AH+AB\ell)+(TBH+ABD)=T(A+B)H+2ABD.
\]
With these real multiplicities, the area still to be removed is $|Y|$ minus the primary area and the desired number $m$ of tags. This gives
\begin{equation}\label{eq:rect-residue}
 R_*=AB+(A+B)(2D+1)T+(H+2D+1)T^2-1-\rho.
\end{equation}
After rounding, let $R$ be the number of additional points to remove from $Y$ to leave exactly $m$ tags. We call these points the private holes. The difference $R-R_*$ is the decrease in primary area due to rounding. In particular $R\ge R_*>0$, since $\rho<AB$. Set $P_x=\lfloor(A-1)/t\rfloor$ and $P_y=\lfloor(B-1)/t\rfloor$. Take $R$ private holes from
\[
 \{at:U_0<a\le P_x\}\times\{at:1\le a\le P_y\}\times[2D,H-1].
\]
The next subsection verifies this set has at least $R$ points. Write $Z_K$ for the chosen set, and remove it together with the primary holes to obtain $|C|=m$. The private vertices of route $(i,b)$ form the set $K_{ib}=Z_K-\mathbf v_{ib}$ in $X$. The first coordinates of the private owners lie outside the first-coordinate intervals of the primary owners. The residues of a private vertex's first two coordinates modulo $t$ determine $i$ and $b$, so the blocks for distinct routes are disjoint.

Use $x$ as the anchor of the loop at $x\in X$. A cycle cover omits this loop whenever its translated tag lies in a hole. If an owner $x$ of $\mathcal R$ in row $i$ serves as a residual anchor at exit $b$, then $x+\mathbf v_{ib}\notin\mathcal R$. Since $x$ belongs to no other primary owner set and to no private block, $x+\mathbf v_{ib}\in C$.

For each pattern in a row, assign distinct residual anchors to as many shared arcs as possible. Assign distinct external anchors of the same pattern to the remaining arcs. The net demand bounds guarantee enough such anchors. The external pools in this row are disjoint and lie outside $X$. A point used twice as a residual anchor serves disjoint intervals of exits, with its loop absent in both. Each cycle cover uses a single row, so reusing an anchor in another row cannot repeat a tag within that cover. Thus the loops and shared arcs in each cycle cover have distinct anchors, and translation by the same $\mathbf v_{ib}$ gives them distinct tags.

If the route passes through $a$ shared owner vertices, its cycle cover contains $a$ shared arcs and $|X|-a-R$ loops. The shared arcs and loops use $m-1-R$ distinct tags. Assign the remaining $R+1$ tags bijectively to the $R+1$ arcs of its private path. Each such arc occurs only on this route, so subtracting $\mathbf v_{ib}$ gives it a fixed anchor. By \cref{eq:main-exposure,lem:exposed-translations}, this gives precisely $t^2$ edge-disjoint perfect matchings.

\subsection{Rectangle counts and the size of the active support}\label{app:rect-support}

\begin{proof}[Proof of \cref{prop:rectangular-factor}]
We bound the number of levels needed for the primary rectangles, count the positions available for private holes, and show that the hub and all owners occupy at most $m/k$ projection vertices. The choice of $T$ gives $T\ge224$ and $128k^3T^7\le m<128k^3(T+1)^7$. Since $AB=24k^2T^4$ and $(1+1/T)^7\le1+7/T+22/T^2$ for $T\ge100$, the height satisfies
\begin{equation}\label{eq:rect-height}
 \frac{16}{3}kT^3-2\le H\le\frac{16}{3}kT^3+39kT^2 .
\end{equation}
In particular $H-3D\ge5kT^3$ and $Q\ge5kT^2$.

First consider the solution of \cref{eq:rect-calibration}. From $S-A=3AD/(H-3D)$ we get $S\le A+6T/5$. Also $AB/Q\le A+7T/5$, since $(A+7T/5)Q-AB=2kT^3/5+(4k-14/5)T^2+7T/5>0$. These bounds give $D(T-1)\{S+(AB-S)/Q\}\le9kT^5/2$. The denominator defining $z$ is positive, since
\[
 \frac{3D}{H}+\frac{D(T-1)}{QH}
 \le\frac{3}{10kT}+\frac1{50k^2T^2}\le\frac1{3T}.
\]
This denominator is at most $BH$. As $BH\ge30k^2T^5$, this yields $z\ge T-3/(20k)>T-1$. To bound $z$ from above, note that both $B/(B-2T)$ and $H/(H-3D-D(T-1)/Q)$ are at most $1+1/(2T)$. By \cref{eq:rect-height}, $AD/H\le3T/8+9/(40T^2)$. Dropping the subtracted term in the numerator now gives
\[
 z\le\left(1+\frac1{2T}\right)^2
       \left(\frac{11T}{8}+\frac9{40T^2}\right)\le\frac{11T}{8}+2 .
\]
Subtracting $11T/8$ from the middle expression gives a remainder that decreases with $T$ and is less than 2 at $T=2$. Since $S\ge A=4kT^2$ and $z\le11T/8+2$, we have $x=S-z>1$. Since $0<x<S<Q$ and $AB/Q>A$, also $y=(AB-x)/Q-z>A-1-z>2$.

Each horizontal level holds $L_h=3kT$ rectangles. For a vertical level,
\[
 \frac{H-D}{2D}
 \ge\frac{16kT}{3}+\frac{16kT}{3(T-1)}-\frac12-\frac{2}{T(T-1)}
 >\frac{16kT}{3}+1.
\]
Taking the floor gives $L_v\ge16kT/3$. Using $n_h\le S$, $n_v\le AB/Q$ and $n_f\le z$ in \cref{eq:rect-level-count} gives
\begin{equation}\label{eq:rect-levels}
 U_0\le\left(\frac{4T}{3}+\frac75\right)
      +\left(\frac{3T}{4}+\frac{101}{80}\right)
      +\left(\frac{11T}{8}+3\right)
 \le\frac{83T}{24}+6.
\end{equation}

For private holes, \cref{eq:rect-residue,eq:rect-height} give
\[
 R_*\le\frac{46}{3}kT^5+(24k^2+39k+1)T^4+10kT^3+T^2
      \le\frac{46}{3}kT^5+65k^2T^4.
\]
Rounding reduces the horizontal area by less than $T(H-2D)$ and the vertical area by less than $2(B-T)D$. After accounting for the shorter rectangle, if present, the area lost from the full type is at most $(B-T)(D+1)$. In all, $0\le R-R_*<TH+3BD+B<15kT^4$, the last inequality following from $43/3+39/T+6/T^2<15$. Consequently
\begin{equation}\label{eq:rect-private}
 0<R\le \frac{46}{3}kT^5+80k^2T^4.
\end{equation}
The floor in $P_x$ gives $P_x=4k(T-1)+\lfloor(4k-1)/(T+1)\rfloor$, and likewise $P_y\ge6k(T-1)$. Thus $P_x-U_0\ge k(13T/24-10)>0$, and $H-2D\ge k(16T^3/3-T^2)$. The number of available private points is at least
\[
 k^3\left(\frac{13T}{24}-10\right)(6T-6)
       \left(\frac{16T^3}{3}-T^2\right)
 \ge k^3\left(\frac{52T^5}{3}-341T^4\right).
\]
Expanding the product gives coefficients $52/3,-4087/12,1533/4,-60$ in degrees $5,4,3,2$, respectively. Subtracting the upper bound in \cref{eq:rect-private} leaves at least $k^3T^4(2T-421)>0$. Since $U_0\le P_x=\lfloor(A-1)/t\rfloor$, the primary levels $t,2t,\ldots,U_0t$ all lie in $[T,A-1]$.

Take $U$ to consist of the hub and every owner. The owners of rectangles on any one primary level occupy at most $t$ first coordinates, with at most $BH$ points per coordinate in the full layers of $X$. The $t^2$ private blocks each have $R$ vertices. Hence $|U|\le 1+tU_0BH+t^2R$. Every projection vertex outside $U$ has only its loop. For $k=1$ we already have $|U|\le m$. For $k\ge2$, using $m\ge ABH$, $H\ge5kT^3$, and \cref{eq:rect-levels,eq:rect-private}, we obtain
\[
 \frac{k|U|}{m}
 \le \frac{83}{96}+\frac{227}{96T}+\frac{3}{2T^2}
   +\left(1+\frac1T\right)^2\left(\frac{23}{180k}+\frac{2}{3T}\right)
   +\frac1{120k^2T^7}<1 .
\]
Indeed, the first three terms sum to less than $9/10$ for $T\ge100$. The fourth is at most $(101/100)^2(23/360+1/150)<3/40$ and the last is less than $1/40$. This proves the support bound.
\end{proof}

For fixed $k$ and $T\to\infty$, \cref{eq:rect-height} gives $H\sim(16k/3)T^3$ and hence $m\sim128k^3T^7$. By \cref{eq:rect-calibration}, $n_h$ and $n_v$ are both asymptotic to $4kT^2$, and $n_f\sim11T/8$. Since rounding reduces the primary area by $O_k(T^4)$, \cref{eq:rect-residue} gives $R\sim(46k/3)T^5$. Summing the owner counts from \cref{lem:rectangle-owners} over all $t$ rows gives contributions of $1/3$, $3/16$, and $11/32$ to the limit of $k|U|/m$ from the horizontal, vertical, and full rectangles, respectively. The shorter rectangle contributes $O_k(T^6)$ owners in total. The $t^2R$ private owners give the last summand in
\[
 \lim_{T\to\infty}\frac{k|U|}{m}
 =\frac13+\frac3{16}+\frac{11}{32}+\frac{23}{192k}
 =\frac{83}{96}+\frac{23}{192k}.
\]
For this choice of $U$ and $k=1$, we have $|U|/m\to63/64$. To fit several disjoint active supports in those parts, we scale the box dimensions with $k$.

\subsection{A planar construction}\label{app:planar}

\begin{proposition}\label{prop:planar-factor}
Let $k\ge1$ and $2\le T<224^2$ be integers. For every integer $m>128k^3T^{7/2}$, there is a balanced tripartite hypergraph of part size $m$ that is $(T+1)$-rigid and has an active support $U$ satisfying $k|U|\le m$.
\end{proposition}

\begin{proof}
Put $t=T+1$, $\gamma_b=\binom b2$, $D=\gamma_T$, $p=32k$, $w=pt$, and $B=(T-1)^2+1$. The planar translations are $\mathbf v_b=(b,\gamma_b)$ for $0\le b\le T$. The functional $\lambda_b(x,y)=2y-(2b-1)x$ takes its unique minimum among the translations at $\mathbf v_b$, since $\lambda_b(\mathbf v_c-\mathbf v_b)=(c-b)^2$. The projection has one row, with $t$ routes through $s$ of the form
\[
 s,S_1,\ldots,S_b,K_b,P_b,\ldots,P_{T-1},s,\qquad 0\le b\le T,
\]
with the same suffix, prefix and private chain conventions as in \cref{app:rectangular}.

Put $R_0=(T^3-T)/2+w-1$. For the box height $h$, we require $h\ge h_0$, where
\begin{equation}\label{eq:planar-height}
 h_0=\max\left\{4D+B,\ 3D+\left\lceil\frac{R_0+D+2B-1}{p-2}\right\rceil\right\}.
\end{equation}
We have $4D+B\le3T^2$ and $p-2\ge30k$. For the second term in the maximum,
\[
 3D+\left\lceil\frac{R_0+D+2B-1}{p-2}\right\rceil
 \le\frac{T^3}{60k}+\frac{19T^2}{12}+\frac{16(T+1)}{15}+1
 \le\frac{T^3}{60k}+3T^2.
\]
Multiplying by $w=32k(T+1)$ gives $wh_0\le(8/15)T^4+145kT^3$. Put $x=\sqrt T$. Since $k\ge1$, the last inequality gives $wh_0/(kT^{7/2})\le 8x/15+145/x$. The function $x\mapsto 8x/15+145/x$ is convex on $[\sqrt2,224]$ and is less than 128 at both endpoints. Consequently $wh_0<128kT^{7/2}\le128k^3T^{7/2}<m$. Write $m-1=wh+\rho$, where $0\le\rho<w$. Then $h\ge h_0$.

Let $X$ consist of $[w]_0\times[h]_0$ and $\rho$ arbitrary points in the next row $[w]_0\times\{h\}$. Use the tag region
\[
 Y=\bigcup_{b=0}^T\bigl([w]_0\times[0,h]+\mathbf v_b\bigr),
 \qquad |Y|=(w+T)(h+1)+(w-1)D .
\]
As $b$ increases, both translation coordinates are nondecreasing. Any point of rectangle $b$ that lies in an earlier rectangle therefore also lies in rectangle $b-1$. Thus the $b$th rectangle adds $(h+1)+(w-1)(\gamma_b-\gamma_{b-1})$ points for $b\ge1$.

\emph{The two sets of primary holes.}
Remove
\[
\begin{aligned}
 Z_1&=\bigl([T,2T-1]\times[D,h-1]\bigr)\cup\bigl([2T,w-1]\times[D,2D-1]\bigr),\\
 Z_2&=[3T,4T-1]\times[3D,3D+B-1].
\end{aligned}
\]
All inverse translates are in the full $w$ by $h$ box. Those of $Z_1$ lie in the strips $x<2T$ or $y<2D$. Those of $Z_2$ lie in $[2T,4T-1]\times[2D,3D+B-1]$. The 2 owner sets are therefore disjoint.

For a point $x$ whose translated tag lies in $Z_1$ at some exit, the exits with $x+\mathbf v_b\in Z_1$ form an interval. Indeed, both coordinates of $x+\mathbf v_b$ are nondecreasing in $b$. At first coordinate $2T$, the translate enters the horizontal arm and the upper bound on its height drops. The interval is proper, since a point whose translate at exit 0 lies in $Z_1$ has first coordinate at least $T$ and height at least $D$, so its translate at exit $T$ misses $Z_1$. Leave a prefix or suffix interval unchanged. Enlarge an internal interval $[a,b]$ to the prefix $[0,b]$ and use the point as a residual anchor on $[0,a-1]$.

For $1\le j\le T$, each of the patterns $[0,j-1]$ and $[j,T]$ then has net external demand
\[
 \sigma_j=h-2D+\gamma_j+(w-2T+j-1)(j-1).
\]
To count suffixes starting at $j$, first consider points with first coordinate $T-j$ before translation. Entry through the vertical arm gives $h-2D+\gamma_j$ points. At the next $j-1$ first coordinates, entry through the lower height boundary gives $j-1=\gamma_j-\gamma_{j-1}$ points each. The horizontal arm contributes a further $(w-2T)(j-1)$. At each of the exits $j-1$ and $j$, exactly $|Z_1|$ owners have translates in $Z_1$. Hence the number of intervals ending at $j-1$ equals the number starting at $j$. After accounting for the residual anchors on the prefix, the two net demands agree.

For a fixed owner, the $T+1$ translates have distinct first coordinates. Since $Z_2$ has only $T$ columns, the owner's interval in $Z_2$ is proper. Enlarge an internal interval $[a,b]$ to the suffix $[a,T]$ and use the point as a residual anchor on $[b+1,T]$. The number of points whose tags lie in $Z_2$ at both exits $0$ and $b$ is $I_b\coloneqq |(Z_2-\mathbf v_0)\cap(Z_2-\mathbf v_b)|=(T-b)(B-\gamma_b)$, since $B\ge D$. Thus $Z_2$ contributes $I_{j-1}-I_j=B-\gamma_j+(T-j+1)(j-1)$ to the prefix demand. Since at every exit $b$ exactly $TB$ points $x\in X$ satisfy $x+\mathbf v_b\in Z_2$, the same count gives the net suffix demand. Together the 2 sets require
\[
 E_j=\sigma_j+I_{j-1}-I_j=h-T+2+(w-T)(j-1)
\]
external anchors for each of the two patterns.

\emph{Anchors on the shared chains.}
For the prefix take the right column and top strip, and for the suffix take the left column and bottom strip, as follows.
\[
\begin{array}{@{}l@{\quad}l@{\quad}l@{}}
 \text{Pattern}&\text{Column}&\text{Strip}\\
 \midrule
 {}[0,j-1]&\{w+T-j\}\times[D-\gamma_{j-1},h+D-\gamma_{j-1}]
         &[T,w-1]\times[h+1+D-\gamma_j,h+D-\gamma_{j-1}]\\
 {}[j,T]&\{-j\}\times[-\gamma_j,h-D+\gamma_{T-j}]
      &[0,w-T-1]\times[-\gamma_j,-\gamma_{j-1}-1]
\end{array}
\]
Also use the following lower intervals on the last $T$ columns, where $0\le a<T$.
\[
 \{w-1-a\}\times[-F_a(j),-F_a(j-1)-1],\qquad
 F_a(b)=\gamma_b-\gamma_{\max\{0,b-a\}} .
\]
Here $F_a(b)$ records the rise in the lower boundary of $Y$. For each column these intervals for different $j$ are disjoint. The bottom strip occupies only earlier columns. All the displayed anchor sets are pairwise disjoint and lie outside $X$.

We check that every allowed translate of these anchors lies in $Y$. For the right-hand column assigned to the prefix $[0,j-1]$, put $q=T-j$. At translated first coordinate $w+q+b$ with $b\le j-1$, the vertical section of $Y$ is $[\gamma_{q+b+1},h+D]$. This section contains every translated anchor in that column because $\gamma_{q+b+1}-\gamma_b$ is nondecreasing in $b$ and $h+D-\gamma_b$ is nonincreasing. For the left-hand column assigned to the suffix $[j,T]$, the section at $b-j$ is $[0,h+\gamma_{b-j}]$. Here $b\ge j$ and $D-\gamma_b\ge \gamma_{T-j}-\gamma_{b-j}$. In the last $T$ columns the lower boundary translated back to column $w-1-a$ is $-F_a(b)$, which is nonincreasing in $b$. Thus the extra lower intervals are valid throughout the suffix, with translated height below $D$. The top strip for the prefix translates into heights $[h+1,h+D]$ at first coordinates at least $T$. The bottom strip for the suffix translates into $[0,w-1]\times[0,D-1]$. An anchor from the right-hand column translates to a first coordinate greater than $w-1$, and one from the left-hand column to a first coordinate less than $T$. Hence all these tags avoid both primary sets and the private holes specified below.

The prefix pool has size $h+1+(w-T)(j-1)=E_j+T-1$. The last $T$ columns add $\sum_{a=0}^{T-1}\min\{j-1,a\}=D-\gamma_{T-j+1}$ suffix anchors. The complete suffix pool has size
\[
 h+1-j(T-j)+(w-T)(j-1)+D-\gamma_{T-j+1}=E_j+\gamma_{j+1}-1\ge E_j.
\]
After the residual anchors have been assigned, each chain has at least as many external anchors available as there are shared arcs still to assign.

\emph{Private holes and completion.}
The primary holes number $Th+(w-3T)D+TB$. Thus $R=|Y|-m-|Z_1|-|Z_2|=(T^3-T)/2+w-\rho-1=R_0-\rho>0$. Choose $R$ points from the union of the three disjoint sets
\[
\begin{aligned}
 &\{3T\}\times[3D+B,h-1],\\
 &\{4T+1\}\times\bigl(\{3D\}\cup[4D+B,h-1]\bigr),\\
 &\{T-2+at:4\le a<p\}\times[3D,h-1].
\end{aligned}
\]
The union has size $(p-2)(h-3D)-(D+2B-1)\ge R_0$ by \cref{eq:planar-height}. Call the chosen set $Z_K$. The inverse sets $K_b=Z_K-\mathbf v_b$ are disjoint by their first-coordinate residues modulo $t$. They lie in $X$ and avoid the two strips containing the $Z_1$ owners. Write a point in the first column as $z=(z_1,z_2)$. If it satisfied $z-\mathbf v_b\in Z_2-\mathbf v_c$, its first coordinate would force $c\ge b$, and hence $z_2\le3D+B-1$. This contradicts the choice of heights. In the second column the same condition forces $b\ge c+2$, giving $3D+1\le z_2\le4D+B-1$, also excluded. The inverse translates of points in all later columns have first coordinate at least $4T+2$.

Take $C=Y\setminus(Z_1\cup Z_2\cup Z_K)$, of size $m$. Each residual anchor is used outside the interval where its translate lies in the original hole set, but inside the enlarged interval where its loop is absent. For each of $Z_1,Z_2,Z_K$, take the union of its translates by $-\mathbf v_b$ over all exits $b$. These three unions are pairwise disjoint, so the residual translate also misses the other holes. The external pools are disjoint and lie outside $X$, so the assignment by pattern in \cref{app:rectangular} applies.

The positive counts $\sigma_j$ ensure that every shared chain is nonempty. Insert the indicated routes and the loops on $X$. Every directed cycle of length at least 2 passes through $s$ and uses a single private path. In each cycle cover, the shared arcs and loops use $m-1-R$ distinct tags. Assign the remaining $R+1$ tags bijectively to the $R+1$ arcs of the private path, which occur only in that cover. For each $b$, the functional $\lambda_b$ attains its minimum over the translations only at $\mathbf v_b$. Hence \cref{lem:exposed-translations} rules out every other perfect matching.

\emph{Active support.}
The hub and all owners form an active support $U$. The union of the two strips for $Z_1$ contains at most $2Th+2D(w-2T)$ owners. The rectangle containing the $Z_2$ owners has $2T(D+B)$ points. Enlarging an owner's interval changes only the exits at which its loop is absent. The private blocks have $tR$ vertices in all. It follows that
\[
 |U|\le 2Th+2D(w-T)+2TB+tR+1.
\]
Here $2Th<m/(16k)$. Using $R\le R_0$, we bound the sum of the remaining terms by
\[
 \frac{T^4}{2}+\left(32k+\frac32\right)T^3+\left(32k-\frac72\right)T^2
       +\left(32k+\frac52\right)T+32k.
\]
Subtract $T^4/2$ from the preceding polynomial and divide the remainder by $T^3$. The result decreases for $T\ge2$ and equals $60k+3/8\le61k$ at $T=2$. Set $x=\sqrt T$. Since $m>128k^3T^{7/2}$,
\[
 \frac{k|U|}{m}
 <\frac1{16}+\frac{1}{128k^2}\left(\frac{x}{2}+\frac{61k}{x}\right)
 \le\frac1{16}+\frac{1}{128k}\left(\frac{x}{2}+\frac{61}{x}\right)<1.
\]
The convex function $x/2+61/x$ is below $113$ at both ends of $[\sqrt2,224]$, and $1/16+113/(128k)\le121/128$. This proves the proposition.
\end{proof}

\subsection{Completion at every part size}

\begin{proof}[Proof of \cref{thm:growth}]
With $k$ and $u$ as above, first suppose $u\le1$. A single perfect matching suffices. If $1<u\le2$, then $m\ge2$. Starting with the two matchings of a bipartite $2m$-cycle, the cyclic extension in \cref{sec:third-part} doubles their number for each of the $k$ added parts, giving $2^{k+1}\ge u^k$ matchings.

For $2<u<224^2$, take $T=\lceil u\rceil-1$. Then $2\le T<224^2$, $T<u\le T+1$, and $m>128k^3T^{7/2}$. By \cref{prop:planar-factor,thm:support-product}, the product of $k$ planar factors has $(T+1)^k\ge u^k$ matchings.

For $u\ge224^2$ we have $m\ge128k^3\,224^7$. Take $T=\lfloor\sqrt u\rfloor$. By \cref{prop:rectangular-factor,thm:support-product}, the product of $k$ rectangular factors has $(T+1)^{2k}>u^k$ matchings. All four constructions have exactly $m$ vertices per part. Since every balanced $r$-partite example is also an $r$-graph on $rm$ vertices, $f_r(rm)\ge g_r(m)$.
\end{proof}

\section{The 9-vertex upper bound}\label{app:finite}

We prove that no $14$-rigid $3$-graph on 9 vertices exists. This is enough for the upper bound in \cref{thm:finite-exact}. If a $d$-rigid $3$-graph existed for $d\ge14$, retaining the edges of any 14 of its perfect matchings would give a $14$-rigid $3$-graph. Throughout this appendix, $H$ denotes a hypothetical $14$-rigid $3$-graph on 9 vertices. Thus $|H|=42$ and every vertex has degree 14.

For a pair $ab$, put $W\coloneqq V(H)\setminus\{a,b\}$ and $X\coloneqq\{x\in W\colon abx\in H\}$. Define the auxiliary graph for this $X$ on all 35 triples of $W$ by joining two triples when they are disjoint and the unique point of $W$ outside their union belongs to $X$. Color each edge by that omitted point. Call the triples of $G\coloneqq H\cap\binom W3$ selected. The subgraph they induce consists of $|X|$ disjoint edges, one of each color, together with isolated vertices. Indeed, a disjoint pair omitting $x\in X$ extends $abx$ to a perfect matching of $H$, and rigidity gives exactly one such extension. Two of these pairs cannot share a triple, because their perfect matchings would then share an edge.

\subsection{Excluding pair codegrees above 4}

Let $X$ be a 5-element set. We shall bound the number of vertices that can be selected in the following bipartite graph. Its two parts are copies of $\binom X2$. Join a left pair to a right pair when they are disjoint, and color the edge by the unique point of $X$ outside their union.

\begin{claim}\label{clm:five-pair-observation}
Select vertices of this graph so that the induced subgraph has maximum degree at most 1 and its edges have distinct colors. If it has 5 edges, exactly 10 vertices are selected. If it has 4 edges, at most 9 vertices are selected.
\end{claim}
\begin{proof}
Write $(P_i,Q_i)$, $i\in[\ell]$, for the matched left and right pairs, and let $x_i$ be their distinct omitted points. Thus $P_i\cap Q_i=\varnothing$, whereas $P_i\cap Q_j\ne\varnothing$ for $i\ne j$. For $x\in X$, let $\alpha_x$ and $\beta_x$ count its occurrences among the $P_i$ and $Q_i$, respectively. Then
\[
 \ell(\ell-1)\le\sum_x\alpha_x\beta_x.
\]
If $\ell=5$, then $\alpha_x+\beta_x=4$ for every $x$, so the right side is at most $20$. Equality holds throughout. Hence $\alpha_x=\beta_x=2$, and $P_i$ and $Q_j$ meet in exactly one point whenever $i\ne j$. In particular the two families of pairs are disjoint. Each is the edge set of a simple $2$-regular graph on 5 vertices, hence a $5$-cycle, and together they partition $\binom X2$. An additional selected left pair would have to intersect every edge of the right $5$-cycle. This is impossible because a $5$-cycle has no vertex cover of size 2. The same applies on the other side. Thus there are exactly 10 selected vertices.

Suppose next that $\ell=4$, and let $x_*$ be the point not among the 4 omitted points. Now $\alpha_x+\beta_x=3$ for $x\ne x_*$ and $\alpha_{x_*}+\beta_{x_*}=4$. Consequently
\[
 12\le\sum_x\alpha_x\beta_x\le4\cdot2+4=12.
\]
Again $|P_i\cap Q_j|=1$ whenever $i\ne j$. Also $\alpha_{x_*}=\beta_{x_*}=2$, and the other 4 $\alpha_x$ belong to $\{1,2\}$. Since $\sum_x\alpha_x=8$, the left pairs form a graph with degree sequence $(2,2,2,1,1)$. This graph is either a path on 5 vertices or the disjoint union of a triangle and an edge.

We rule out the second case. A right pair meets exactly three of the four left pairs and is distinct from all four. It would therefore join a triangle vertex to an endpoint of the isolated edge. Every right pair would then meet that isolated edge, which must have a disjoint partner.

We may therefore label $X=[5]$ so that the left pairs are $12,23,34,45$. A right pair meets exactly three of these pairs and is distinct from all four. The possibilities are $13,14,25,35$, so these are all the right pairs. An additional selected left pair must meet all 4 right pairs. The only 2-element vertex cover of their path is $15$. Similarly the only possible additional right pair is $24$. These 2 pairs are disjoint, so selecting both creates a fifth edge in the induced graph. At most one can be added, proving the upper bound 9.
\end{proof}

\begin{lemma}\label{lem:four-codegree}
Every pair of vertices in a $14$-rigid $3$-graph on nine vertices has codegree at most four.
\end{lemma}
\begin{proof}
Suppose a pair $a,b$ has codegree $k\ge5$. With $W,X,G$ as above, $|W|=7$, $|X|=k\le7$, and $|G|=42-(14+14-k)=14+k$. In the auxiliary graph for this $X$, the subgraph induced by $G$ consists of $k$ disjoint edges, one of each color in $X$, and isolated vertices.

If $k=7$, the full graph on the $35$ triples is $4$-regular. Its selected set $G$ has size $21$ and induces 7 edges. Thus the cut from $G$ to its complement has size $4\cdot21-2\cdot7=70$, but the 14 vertices outside $G$ are incident with at most $4\cdot14=56$ cut edges.

If $k=6$, write $W=X\cup\{y\}$. The auxiliary graph for this $X$ is bipartite. Triples containing $y$ are identified with pairs of $X$, and a triple contained in $X$ is identified with its 3-element complement in $X$. Adjacency is precisely pair containment in a triple. The left and right degrees are 4 and 3, respectively. Let $p$ left vertices and $q=20-p$ right vertices be selected. The 6 induced matching edges imply $p\ge6$. Counting edges from the selected left vertices to the unselected right vertices gives $4p-6\le3(20-q)=3p$, so $p=6$. Equality means that each of the 6 unselected right triples has all 3 of its pairs among the 6 selected left pairs. The graph on $X$ formed by the six selected pairs would therefore contain at least six triangles. But each triangle uses 3 pairs of adjacent edges, and each pair of edges belongs to at most one triangle. Hence a graph with 6 edges has at most $\binom62/3=5$ triangles, a contradiction.

It remains to consider $k=5$. Write $W=X\cup\{u,v\}$. The auxiliary graph for this $X$ has two components. The first has triples $uvx$ on one side and triples contained in $X$ on the other. Identifying each triple in $X$ with its complementary pair gives the vertex--edge incidence graph of $K_5$. The second has triples $uP$ and $vQ$, with $P,Q\in\binom X2$, and joins them when $P\cap Q=\varnothing$. It is the graph of \cref{clm:five-pair-observation}.

Let $s$ of the 5 induced matching edges lie in the incidence component. Let $A\subseteq X$ be the set of vertices of $K_5$ selected on the left of its incidence graph. Each selected vertex on the other side represents an edge of $K_5$ with at most one endpoint in $A$. Exactly $s$ of these edges have an endpoint in $A$, and the others lie in $X\setminus A$. Thus the number $g_1$ of selected vertices in this component satisfies
\[
 g_1\le |A|+\binom{5-|A|}{2}+s.
\]
If a selected incidence edge joins $x\in A$ to the edge $\{x,y\}$ of $K_5$, its color is $y\in X\setminus A$. Since the $s$ colors are distinct and the induced edges form a matching, $s\le\min\{|A|,5-|A|\}\le2$.

The other component is a $3$-regular graph on 20 vertices. If it contains $g_2$ selected vertices and $\ell=5-s$ induced edges, counting the edges from the $g_2$ selected vertices to the $20-g_2$ unselected vertices gives
\[
 3g_2-2\ell\le3(20-g_2),
 \quad\text{and hence}\quad g_2\le10+\left\lfloor\frac\ell3\right\rfloor.
\]
If $s=2$, then $|A|\in\{2,3\}$, so $g_1\le7$ and $g_2\le11$, contradicting $g_1+g_2=|G|=19$. If $s=1$, then $g_1\le8$, while \cref{clm:five-pair-observation} gives $g_2\le9$, again a contradiction.

Thus $s=0$. The same claim gives $g_2=10$, and hence $g_1=9$. The bound for $g_1$ forces $A=\varnothing$. Consequently $G$ consists of 9 of the 10 triples contained in $X$, together with 10 triples of the form $uP$ or $vQ$. These ten triples are the endpoints of the five induced edges, so they are all used in the 5 perfect matchings containing $abx$, $x\in X$.

There are 9 remaining perfect matchings of $H$. Each contains one of the 9 triples of $X$, one edge containing $a$ but not $b$, and one edge containing $b$ but not $a$. There is at most one edge of the form $axy$ with $x,y\in X$, because any such edge must be paired with the unique possible edge $buv$ in its perfect matching. Similarly there is at most one edge $bxy$ with $x,y\in X$. There is also at most one edge $auv$ and at most one edge $buv$. Hence at least 7 of the 9 edges containing $a$ but not $b$ have the form $aux$ or $avx$ for some $x\in X$.

After interchanging $u$ and $v$ if necessary, we may assume that at least 4 edges have the form $aux$. They belong to distinct perfect matchings. Each of these matchings pairs its edge $aux$ with an edge $bvy$, and the latter edges are distinct by rigidity.

Fix one such $aux$. There are at least four choices of $y$ with $bvy\in H$. At most one equals $x$, and at most one other choice makes $X\setminus\{x,y\}$ the unique triple of $X$ missing from $G$. Each remaining $y$ gives a perfect matching $\{aux,bvy,X\setminus\{x,y\}\}$. At least 2 of them therefore contain $aux$, contradicting rigidity.
\end{proof}

\subsection{The triples avoiding a pair of codegree four}

\begin{lemma}\label{lem:four-neighborhood}
Suppose that $H$ is a $14$-rigid $3$-graph on 9 vertices and that every pair has codegree at most 4. There is a pair $ab$ of codegree 4 for which one can label the other 7 vertices $0,1,\ldots,6$ so that $\{x\colon abx\in H\}=\{0,1,2,3\}$. With this labeling, the edges avoiding $a,b$ are
\[
\begin{aligned}
 \{&012,013,014,015,023,024,026,034,035,046,056,\\
     &123,125,136,146,236,245,345\}.
\end{aligned}
\]
Here, for example, $012$ denotes $\{0,1,2\}$.
\end{lemma}
\begin{proof}
Since the sum of the 36 pair codegrees is $3|H|=126>3\cdot36$, some pair $ab$ has codegree 4. With $W,X,G$ as defined above for this pair, we have $|X|=4$ and $|G|=18$. Put $Y\coloneqq W\setminus X$. In the auxiliary graph for $X$, the subgraph induced by $G$ has four disjoint edges, one of each color in $X$.

For $y\in Y$, there is no edge $aby$. At most 4 edges contain $a,y$, and at most 4 contain $b,y$, because pair codegrees are at most 4. Hence
\begin{equation}
 \deg_G(y)\ge14-4-4=6\qquad(y\in Y).
 \label{eq:four-G-minimum-degrees}
\end{equation}

Write $Y=\{y_1,y_2,y_3\}$. The auxiliary graph on the 35 triples of $W$ has four components. One is a star $K_{1,4}$, with center $Y$ and leaves $\binom X3$. Index the three remaining components by $i\in[3]$. Component $i$ has left vertices $(Y\setminus\{y_i\})\cup\{x\}$, $x\in X$, and right vertices $\{y_i\}\cup P$, $P\in\binom X2$. Identifying each right vertex $\{y_i\}\cup P$ with the pair $X\setminus P$ gives the vertex--edge incidence graph of $K_4$. Its incidence edge joining a vertex $x$ to a pair $\{x,x'\}$ has color $x'$.

Let $a_i,b_i,s_i$ count the selected left vertices, selected right vertices, and induced matching edges in the $i$th incidence component. Each pair labeling a selected right vertex contains at most one selected left label. Exactly $s_i$ pairs contain one, and all the others use two unselected labels. Consequently
\begin{equation}
 b_i\le\binom{4-a_i}{2}+s_i,
 \qquad s_i\le\min\{a_i,4-a_i\}\le2.
 \label{eq:four-component-bound}
\end{equation}
The $s_i$ induced edges have distinct left endpoints among the $a_i$ selected left vertices. Their colors are distinct points among the other $4-a_i$ vertices of $X$. For $s_i=0,1,2$ the maximum numbers $a_i+b_i$ of selected vertices are $6,5,5$, respectively. Attaining these maxima forces $(a_i,b_i)=(0,6),(1,4),(2,3)$, respectively. If $s_i=2$, then always $a_i=2$ and $b_i\in\{2,3\}$.

Let $\varepsilon$ indicate whether the star center $Y$ is selected, and let $\ell$ count its selected leaves. Directly from the component descriptions,
\begin{equation}
 \deg_G(y_i)=\varepsilon+b_i+a_j+a_k,
 \qquad \{i,j,k\}=\{1,2,3\}.
 \label{eq:four-Y-degrees}
\end{equation}
We use this identity to determine how many vertices of $G$ lie on each side of each incidence component.

If the star supplies one of the 4 induced matching edges, then $\varepsilon=\ell=1$. Up to reordering the three incidence components, $(s_1,s_2,s_3)$ is $(2,1,0)$ or $(1,1,1)$. The latter permits at most $2+5+5+5=17$ selected vertices. In the former, 18 selected vertices force every component to attain its maximum, giving $(a_i,b_i)=(2,3),(1,4),(0,6)$. At $y_1$, \cref{eq:four-Y-degrees} gives $\deg_G(y_1)=1+3+1+0=5$, contrary to \cref{eq:four-G-minimum-degrees}.

Thus the star contains no induced matching edge. If its center were selected, it would be the only selected vertex of the star. The three incidence components together contain at most $5+5+6=16$ selected vertices, giving fewer than 18 in all. Hence $\varepsilon=0$ and $0\le\ell\le4$. The 4 matching edges are distributed among the incidence components as $(2,1,1)$ or $(2,2,0)$.

First suppose $(s_1,s_2,s_3)=(2,1,1)$. Each incidence component contains at most 5 selected vertices, so $\ell\in\{3,4\}$. If $\ell=3$, all three contain 5 selected vertices, and \cref{eq:four-Y-degrees} gives $\deg_G(y_1)=3+1+1=5$. If $\ell=4$ and the first component contains 4 selected vertices, the other two contain 5 each, giving $\deg_G(y_1)=2+1+1=4$. Both cases contradict \cref{eq:four-G-minimum-degrees}. Otherwise, after interchanging components 2 and 3 if necessary, the numbers of selected vertices are $5,5,4$. The corresponding $(a_i,b_i)$ are $(2,3),(1,4),(a,b)$ with $a+b=4$. Applying \cref{eq:four-Y-degrees} at $y_1$ and $y_3$ and using the lower bound in \cref{eq:four-G-minimum-degrees} gives $3+1+a\ge6$ and $b+2+1\ge6$. Thus $a\ge2$ and $b\ge3$, contradicting $a+b=4$.

The remaining distribution is $(2,2,0)$. Number the components so that $s_3=0$. The first two have $a_1=a_2=2$ and at most 5 selected vertices each, so the third contains at least 4 selected vertices. By \cref{eq:four-component-bound}, the possible counts on its two sides satisfy $a_3=0$, $(a_3,b_3)=(1,3)$, or $(a_3,b_3)=(4,0)$. The first gives $\deg_G(y_1)\le3+2+0=5$. The last gives $\deg_G(y_3)=0+2+2=4$. Therefore $(a_3,b_3)=(1,3)$. Counting the 18 selected vertices now forces
\begin{equation}
 \ell=4,\qquad
 (a_1,b_1)=(a_2,b_2)=(2,3),\qquad
 (a_3,b_3)=(1,3).
 \label{eq:four-forced-profile}
\end{equation}

Let $P,Q\subseteq X$ consist of the labels of the selected left vertices in the first and second components, respectively. By \cref{eq:four-forced-profile}, each has size 2. The colors of the two matching edges in the first component are the two points of $X\setminus P$, and those in the second are the two points of $X\setminus Q$. Since all 4 colors are distinct, $P$ and $Q$ partition $X$. The 3 selected triples of the form $y_1xx'$ in the first component are $\{y_1\}\cup P$ and the 2 triples obtained by adjoining $y_1$ to the edges of a perfect matching between $P$ and $Q$. The second component likewise gives $\{y_2\}\cup Q$ and 2 triples from a perfect matching between $P$ and $Q$.

Write $x_0$ for the label of the selected left vertex in the third component. This component contributes $y_1y_2x_0$ and the three triples $y_3x_0x$ with $x\in X\setminus\{x_0\}$. The 2 perfect matchings between $P$ and $Q$ must be different. Indeed, if they agree, take their common pair $x_0x$. The pair $x_0x$ then lies in 2 triples of $\binom X3$, in $y_3x_0x$, and in both $y_1x_0x$ and $y_2x_0x$, giving codegree at least 5.

Interchanging the first two incidence components if necessary, assume $x_0\in P$. Relabel so that $P=\{0,1\}$, $Q=\{2,3\}$, $x_0=0$, $y_1=5$, $y_2=6$, and $y_3=4$. Choose the labels $2,3$ in $Q$ so that the first matching is $\{03,12\}$. The second is then $\{02,13\}$. Applying this labeling to the triples described above gives the 18 edges in the statement.
\end{proof}

\subsection{The seven-vertex configuration has no rigid extension}

\begin{lemma}\label{lem:four-normal-form-extension}
There is no $14$-rigid $3$-graph $H$ on $\{a,b,0,1,\ldots,6\}$ with maximum pair codegree at most 4 for which $\{x\colon abx\in H\}=\{0,1,2,3\}$ and $G\coloneqq H[\{0,1,\ldots,6\}]$ has the 18 edges listed in \cref{lem:four-neighborhood}.
\end{lemma}
\begin{proof}
Let $\mathcal A$ consist of the pairs $S\subseteq\{0,\ldots,6\}$ with $aS\in H$, and define $\mathcal B$ likewise for $b$. Both families have size 10, since $H$ is $14$-regular and precisely 4 edges contain both $a$ and $b$.

The perfect matchings through $ab1$, $ab2$, $ab3$, and $ab0$ use, in that order, the following pairs of edges of $G$
\[
 \{026,345\},\quad \{035,146\},\quad
 \{046,125\},\quad \{136,245\}.
\]
These 8 edges of $G$ cannot occur in any further perfect matching. Every remaining perfect matching has the form $\{aS,bT,W\setminus(S\cup T)\}$, where $S\in\mathcal A$, $T\in\mathcal B$, $S\cap T=\varnothing$, and $W\setminus(S\cup T)\in G$.

The pairs $01,02,03,04$ already have codegree 4 in $G$, while $05$ and $06$ have codegree 3. Also $\deg_G(0)=11$, and $ab0\in H$. Consequently $05$ and $06$ each belong to exactly one of $\mathcal A,\mathcal B$, and neither family contains any other pair through $0$. Interchanging $a,b$ if necessary, assume $06\in\mathcal A$.

For $S=06$, the candidate pairs for $T$ are $12,13,34,45$. The first three give the already used triples $345,245,125$, respectively. Thus $45\in\mathcal B$, and the matching containing $a06$ is $\{a06,b45,123\}$. Similarly, the possible partners of $05$ are $14,23,24,46$, with complementary triples $236,146,136,123$, respectively. The last three triples have already been used, so $05$ must be paired with $14$ in the opposite family.

If $05\in\mathcal A$, then $\mathcal B$ contains no pair through $0$. The unique perfect matching containing $a06$ uses $b45$. If any of $12,13,34$ belonged to $\mathcal B$, its complementary triple with $06$ would give another perfect matching containing $a06$. The same argument for $a05$, whose matching uses $b14$, excludes $23,24,46$ from $\mathcal B$. These are 6 distinct pairs among the 15 pairs on $\{1,\ldots,6\}$, so $|\mathcal B|\le15-6=9$, contradicting $|\mathcal B|=10$.

If $05\in\mathcal B$, then $06$ is the only pair of $\mathcal A$ through $0$. The matching containing $b05$ uses $a14$, so $23,24,46\notin\mathcal A$. The unique perfect matching through $b45$ uses $a06$, so $13,16,26,36\notin\mathcal A$, because their complementary triples with $45$ are $026,023,013,012$, all in $G$. These are 7 distinct pairs on $\{1,\ldots,6\}$, whence $|\mathcal A|\le1+(15-7)=9$, contradicting $|\mathcal A|=10$.
\end{proof}

\begin{proof}[Proof of \cref{thm:finite-exact} (upper bound for $f_3(9)$)]
By \cref{lem:four-codegree,lem:four-neighborhood}, any $14$-rigid $3$-graph on 9 vertices would contain the configuration excluded by \cref{lem:four-normal-form-extension}. Hence $f_3(9)\le13$. The constructions and the tripartite upper bound in \cref{sec:finite} complete the proof.
\end{proof}



\bibliographystyle{abbrv}
\bibliography{main}
\end{document}
